\documentclass[11pt]{article}

\usepackage[T1]{fontenc}
\usepackage[utf8]{inputenc}
\usepackage[english]{babel}
\usepackage[margin=1in]{geometry}
\usepackage{lmodern}
\usepackage{microtype}
\usepackage{mathtools}
\usepackage{amsmath,amssymb,amsthm}
\usepackage{array}
\usepackage{booktabs}
\usepackage{tabularx}
\usepackage{float}
\usepackage{graphicx}
\usepackage{pdflscape}
\usepackage{caption}
\usepackage{enumitem}
\usepackage{fancyhdr}
\usepackage[numbers]{natbib}
\usepackage{doi}
\usepackage{xurl}
\usepackage{hyperref}
\usepackage{cleveref}
\usepackage{orcidlink}
\usepackage{titling}



\newcommand{\OldLayer}{H2}
\newcommand{\NewLayer}{H6}

\newcommand{\PaperNine}{Paper~9}
\newcommand{\PaperEight}{Paper~8}
\newcommand{\AcceptedTestCount}{three}
\newcommand{\FinalVerdict}{\textsc{H6\_NEW\_LAYER\_STRICT\_THEORY\_EXTENSION\_WITH\_OLD\_LAYER\_REAL\_PROVISIONAL}}
\newcommand{\OldLayerStatus}{real-provisional non-factorization boundary}

\newcommand{\SlotReadoutAgreement}{\texttt{slot.h6\_pica\_new\_layer\_protocol\_indexed\_readout\_agreement\_family}}
\newcommand{\SlotQuadraticWitness}{\texttt{slot.h6\_pica\_new\_layer\_quadratic\_witness\_family}}
\newcommand{\SlotTransitionWitness}{\texttt{slot.h6\_pica\_new\_layer\_transition\_witness\_family}}

\newcommand{\TestFamilyCompletion}{\texttt{test.h6\_pica.strict\_family.completion\_family\_surface}}
\newcommand{\TestFamilyLifted}{\texttt{test.h6\_pica.strict\_family.lifted\_package\_surface}}
\newcommand{\TestFamilyHigher}{\texttt{test.h6\_pica.strict\_family.protocol\_indexed\_higher\_witness\_surface}}
 
\newtheorem{theorem}{Theorem}[section]
\newtheorem{lemma}[theorem]{Lemma}
\newtheorem{proposition}[theorem]{Proposition}
\newtheorem{corollary}[theorem]{Corollary}
\theoremstyle{definition}
\newtheorem{definition}[theorem]{Definition}
\newtheorem{remark}[theorem]{Remark}
\floatstyle{ruled}
\newfloat{algorithm}{tbp}{loa}
\floatname{algorithm}{Algorithm}

\urlstyle{same}
\captionsetup{
  font=small,
  labelfont=bf,
  labelsep=period
}
\captionsetup[table]{skip=6pt,position=top}
\captionsetup[figure]{skip=6pt}
\captionsetup[algorithm]{skip=6pt,position=top}
\crefname{algorithm}{algorithm}{algorithms}
\Crefname{algorithm}{Algorithm}{Algorithms}
\crefname{theorem}{theorem}{theorems}
\Crefname{theorem}{Theorem}{Theorems}
\crefname{lemma}{lemma}{lemmas}
\Crefname{lemma}{Lemma}{Lemmas}
\crefname{proposition}{proposition}{propositions}
\Crefname{proposition}{Proposition}{Propositions}
\crefname{corollary}{corollary}{corollaries}
\Crefname{corollary}{Corollary}{Corollaries}
\crefname{definition}{definition}{definitions}
\Crefname{definition}{Definition}{Definitions}
\crefname{remark}{remark}{remarks}
\Crefname{remark}{Remark}{Remarks}
\setlist[enumerate]{leftmargin=*, itemsep=2pt, topsep=4pt}
\setlength{\droptitle}{-3.2em}
\pretitle{\begin{center}\LARGE\bfseries}
\posttitle{\par\end{center}\vspace{0.5em}}
\preauthor{\begin{center}\normalsize}
\postauthor{\par\end{center}\vspace{-0.4em}}
\predate{\begin{center}\small}
\postdate{\par\end{center}\vspace{0.6em}}
\setlength{\emergencystretch}{2em}
\clubpenalty=10000
\widowpenalty=10000
\displaywidowpenalty=10000
\interfootnotelinepenalty=10000
\allowdisplaybreaks[2]
\fancypagestyle{firstpage}{\fancyhf{}
  \fancyfoot[C]{\scriptsize Automorph Inc., Wilmington, DE, USA \quad \texttt{ioannis@automorph.io} \quad \textcopyright\ 2026 \quad CC-BY 4.0}
  \renewcommand{\headrulewidth}{0pt}
  \renewcommand{\footrulewidth}{0pt}
}

\title{Accepted New-Layer Strict Tests on Fixed Support in Six Birds Theory}
\author{Ioannis Tsiokos\,\orcidlink{0009-0009-7659-5964}}
\date{\today}
\hypersetup{
  pdftitle={Accepted New-Layer Strict Tests on Fixed Support in Six Birds Theory},
  pdfauthor={Ioannis Tsiokos},
  hidelinks
}

\begin{document}

\maketitle
\thispagestyle{firstpage}

\begin{abstract}
Scoped event admissibility on fixed support leaves a specific gap for strict
tests: the protocol-indexed flagship-seed bridge family realizes the
quadratic-levelset and transition-stability slots, but the
protocol-indexed readout-agreement slot remains unrealized, so no accepted
strict test family is available at the bridge-native event level.  This paper
asks whether strict-test acceptance is admissible on fixed support without
repairing that bridge level.  The method is a layer-local, audit-based PICA
route: promotion, internalization of protocol coordinates, carrier
exactification, and audited source admission produce a \NewLayer{} carrier
substrate above the bridge-native \OldLayer{} level.  On the resulting
\NewLayer{} carriers, the readout-agreement, quadratic-witness, and
transition-witness slots are strict-realized natively; the three-slot strict
domain closes; and \AcceptedTestCount{} accepted scoped strict test families
survive the required control panel.  The panel covers null-baseline,
memory-only, artifact, boundary, selected-lift, visibility, and claim-boundary
checks.  The bridge-native \OldLayer{} level is
retained as the \OldLayerStatus{}: descent from new-layer evidence to
old-layer compatibility preserves non-factorization without inducing
old-layer strict realization.  The final classification is a strict theory
extension at the new layer with old-layer real-provisional compatibility
intact.  The result is bounded: it does not assert event-package construction,
shared-event identity, proto-effect or effect-algebra structure, Born
valuation, composition readiness, faithful public compression, a lawful public
carrier, or canonical selected lifts.
\end{abstract}

\medskip
\noindent\textbf{Keywords.} Six Birds Theory; fixed support; strict theory extension; carrier exactification; protocol-indexed event admissibility; layer localization; non-factorization; claim-boundary discipline.

\section{Introduction}\label{sec:intro}

The fixed-support Six Birds series has progressed from same-support event
admissibility and obstruction
\cite{Tsiokos2026LocallyBoolean,Tsiokos2026HolonomyMemory,Tsiokos2026MarkingErasureRecombination,Tsiokos2026RecombinationWitnesses}
through promoted-package classification and carrier exactification
\cite{Tsiokos2026PromotionCriteria,Tsiokos2026InternalStructure,Tsiokos2026CarrierExactification}
to scoped event admissibility \cite{Tsiokos2026CarrierToEvent}.  The immediate
predecessor, \PaperEight{} \cite{Tsiokos2026CarrierToEvent}, ends with a
three-slot proto-family domain at the bridge-native event level: the
quadratic-levelset and transition-stability slots are realized as scoped event
positives, while the protocol-indexed readout-agreement slot remains
unrealized.  Consequently, no accepted strict test family is available at that
bridge level.  The question of \PaperNine{} is whether strict-test acceptance
is admissible at all on fixed support, and, if so, at which layer.

\PaperNine{} answers by separating the layer of strict-test admissibility from
the inherited bridge-native layer.  The paper does not attempt a direct repair
of the bridge level.  Instead, it constructs a new-layer carrier substrate by
the PICA route: the protocol-indexed carrier-audit route of
\Cref{sec:methods}, combining promotion, internalization of protocol
coordinates, carrier exactification, and audited source admission.  In the
paper-local notation, \OldLayer{} denotes the bridge-native event level
inherited from \PaperEight{}, and \NewLayer{} denotes the new-layer carrier
substrate constructed here.

The principal result is that the \NewLayer{} carriers admit a closed
three-slot strict domain with \AcceptedTestCount{} accepted strict test
families that survive the required control panel.  The bridge-native
\OldLayer{} is retained as the \OldLayerStatus{}: descent from \NewLayer{}
evidence to \OldLayer{} compatibility preserves non-factorization without
inducing old-layer strict realization.  The final classification is therefore a
strict theory extension at the new layer with old-layer real-provisional
compatibility intact; the formal verdict and its legality constraints are
recorded in \Cref{sec:final-gate}.

The result is bounded.  \PaperNine{} does not assert event-package
construction, shared-event identity, proto-effect or effect-algebra structure,
Born valuation, composition readiness, faithful public compression, a lawful
public carrier, or canonical selected lifts.  These nonclaims are inherited and
then enforced by the present paper's claim-boundary discipline; the compact
inheritance statement appears in \Cref{subsec:inherited-nonclaims}, and the
complete ledger appears in \Cref{subsec:nonclaims-ledger}.

\Cref{sec:background,sec:problem-setup} record the inherited state and the
layer-local setup.  \Cref{sec:methods,sec:primary-results,sec:secondary-results}
construct the \NewLayer{} carriers and report the strict-realization, descent,
and control-survival audits.  The remaining sections give the
layer-localization analysis, formal schemas, final gate, discussion, and
conclusion.
 \section{Starting point: scoped event admissibility after \PaperEight{}}\label{sec:background}

\subsection{Earlier papers in the fixed-support series}\label{subsec:earlier-papers-chain}

The fixed-support sequence begins with
\cite{Tsiokos2026LocallyBoolean}, which establishes event packages and the
same-support comparison discipline.  \cite{Tsiokos2026HolonomyMemory}
introduces the current and predictive quotients on single histories.
\cite{Tsiokos2026MarkingErasureRecombination} extends the setting to
branchwise and recombination quotients on branch assemblages.
\cite{Tsiokos2026RecombinationWitnesses} isolates the control-stable Route A
flagship family with a Lean-formal kernel.  \cite{Tsiokos2026PromotionCriteria}
then supplies the promotion criteria and identifies the intrinsic
promoted-package survivor.

The next two papers analyze that survivor.  \cite{Tsiokos2026InternalStructure}
classifies the surviving promoted package and retains the invariant and
quadratic families used downstream.  \cite{Tsiokos2026CarrierExactification}
performs the seven-criterion carrier-exactification audit and obtains the
lawful same-branch witness with its public-shadow split.  The immediate
predecessor, \PaperEight{} \cite{Tsiokos2026CarrierToEvent}, closes this chain
at scoped event admissibility and is treated in
\Cref{subsec:paper8-inherited-contract}.

\subsection{\PaperEight{} inherited contract}\label{subsec:paper8-inherited-contract}

\PaperEight{} raises the carrier atlas of
\cite{Tsiokos2026CarrierExactification} to the event level by introducing
event bridge records and admissibility audits
\cite{Tsiokos2026CarrierToEvent}.  Its principal classification is that the
protocol-indexed flagship-seed bridge family is event-level admissible.  The two
scoped event positives, the quadratic-levelset event and the
transition-stability test, pass the event gate as positively claimed scoped
events.  The same bridge family is not a closed test family: its proto-family
domain has three required slots, and only two are realized at the bridge level.

\PaperNine{} inherits this state.  The bridge-native event level
becomes the \OldLayer{} of the present paper; the two scoped event positives
are retained as inherited evidence; and the unrealized readout-agreement slot is
the gap addressed below.  It also inherits the test-family discipline of
mutual exclusivity, exhaustivity, readout-category coverage, and family
structure used in \Cref{sec:primary-results}; the bridge-native source policy
that \Cref{sec:methods} replaces with a PICA-bound policy at the new layer; and
the explicit nonclaim record, including no event package, no proto-effect
structure, and no faithful public compression.  The readout-agreement gap is
isolated in \Cref{subsec:unrealized-readout-slot}, the inherited nonclaims are
recorded in \Cref{subsec:inherited-nonclaims}, and
\Cref{subsec:inherited-state-table} summarizes the inherited state.

\subsection{Unrealized readout-agreement slot}\label{subsec:unrealized-readout-slot}

At the bridge level of \PaperEight{}, the proto-family domain requires three
slots under the test-family discipline of \cite{Tsiokos2026CarrierToEvent}: a
quadratic-levelset slot, a transition-stability slot, and a protocol-indexed
readout-agreement slot.  The first two are realized as scoped event positives
at that level.  The readout-agreement slot is not realized there because the
bridge-native source policy admits the relevant committed-state evidence only
for provisional inspection, which is insufficient for strict realization under
the threshold-and-witness discipline of \cite{Tsiokos2026FoundationsII}.  This
is the structural opening addressed by \PaperNine{} through a layer change
rather than direct repair: \Cref{sec:methods} describes the new-layer route, and
\Cref{subsec:source-backed-threshold-evidence,subsec:new-layer-readout-slot-gate}
record the new-layer threshold evidence and readout-slot gate in
\Cref{sec:primary-results}.

\subsection{Inherited nonclaims}\label{subsec:inherited-nonclaims}

\PaperEight{}'s scoped event admissibility is bounded by an explicit nonclaim
record \cite{Tsiokos2026CarrierToEvent}, in the foundations claim-discipline
sense of \cite{Tsiokos2026FoundationsII}.  It does not assert event-package
construction, a proto-effect fragment, an effect algebra, a Born valuation,
composition readiness, faithful public compression, a lawful public carrier, or
a canonical selected lift.  These nonclaims survive intact into \PaperNine{}:
they are not reopened by the new-layer construction, strict-test acceptance, or
descent-and-control survival.  The claim-boundary discipline of
\Cref{subsec:pica-claim-boundary-report} preserves the same nonclaim record at
both layers, and the claim-boundary control of
\Cref{subsec:protocol-honesty-claim-boundary} verifies that the accepted
families respect it.

\subsection{Inherited fixed-support state}\label{subsec:inherited-state-table}

\Cref{tab:inherited-state} summarizes the fixed-support state inherited from
\PaperEight{} and carried forward into \PaperNine{}.

\begin{table}[tbp]
  \centering
  \caption{Inherited fixed-support state at the close of \PaperEight{}: status, scope, and disposition carried forward into \PaperNine{}.}
  \label{tab:inherited-state}
  \small
  \begin{tabularx}{\textwidth}{@{}>{\raggedright\arraybackslash}p{0.24\textwidth}>{\raggedright\arraybackslash}X>{\raggedright\arraybackslash}p{0.20\textwidth}>{\raggedright\arraybackslash}p{0.22\textwidth}@{}}
    \toprule
    Inherited element & Status at \PaperEight{} & Scope & Carried forward as \\
    \midrule
    Quadratic-levelset scoped event
      & scoped event positive
      & bridge-native; scoped protocol-indexed
      & inherited evidence at \OldLayer{} \\
    Transition-stability scoped test
      & scoped event positive; pre-test
      & bridge-native; scoped protocol-indexed
      & inherited evidence at \OldLayer{} \\
    Old-layer readout-agreement slot
      & unrealized at the bridge level
      & bridge-native; protocol-indexed
      & gap addressed by layer change \\
    Accepted strict tests at the bridge level
      & none asserted
      & bridge-native
      & old-layer boundary condition \\
    Public shadow / public-lawful carrier
      & nonfaithful public shadow; no lawful public carrier asserted
      & public shadow of the bridge-native carrier
      & nonclaim record \\
    Selected lifts (canonical-lift status)
      & selected lifts nonunique; no canonical selected lift asserted
      & level-crossing records
      & nonclaim record \\
    Higher structures (event package, proto-effect, effect algebra, Born valuation, composition)
      & not asserted
      & beyond \PaperEight{}'s scoped event result
      & nonclaim record \\
    \bottomrule
  \end{tabularx}
\end{table}
 
With the inherited state fixed, \Cref{sec:problem-setup} formulates the
layer-local strict-test problem that the rest of the paper addresses.
 \section{Layer-local strict-test admissibility and the \OldLayer{} baseline}\label{sec:problem-setup}

\subsection{Old-layer bridge baseline}\label{subsec:h2-bridge-baseline}

At the close of \PaperEight{}, the bridge-native \OldLayer{} event level carries
a three-slot proto-family domain
\cite{Tsiokos2026CarrierToEvent}.  Its quadratic-levelset slot and
transition-stability slot are realized as scoped event positives at that level.
Its protocol-indexed readout-agreement slot remains unrealized.  This missing
readout-agreement slot is the local gap inherited by \PaperNine{}: the old layer
contains genuine scoped event structure, but it does not support an accepted
strict test family.

Throughout this paper, \OldLayer{} is retained as the \OldLayerStatus{}.  This
means that the old layer is neither falsified nor strict-realized.  Its scoped
event positives remain valid, its readout-agreement slot remains unrealized,
and it remains the inherited target of any descent from a higher-layer
construction under the foundations layer discipline
\cite{Tsiokos2026FoundationsII}.  \PaperNine{} therefore does not attempt a
direct repair of the bridge-native level.  Its question is whether a
\NewLayer{} construction can supply strict realization while preserving this
boundary.  \Cref{subsec:new-layer-carrier-native-slots} introduces the new-layer
carrier and native slots, \Cref{subsec:non-factorization-baseline} records the
non-factorization requirement, and \Cref{subsec:old-layer-status-table} records
the old-layer status ledger.

\subsection{New-layer carrier and native strict slots}\label{subsec:new-layer-carrier-native-slots}

A new-layer carrier is a finite description with its own native state space,
transport law, and readout law, constructed above the bridge-native event level
of \PaperEight{} by the methodology of \Cref{sec:methods}.  Conceptually, it is
the substrate on which later strict-realization claims may be made.  The state
space supplies the carrier's content domain, the transport law specifies its
dynamics, and the readout law specifies the observable content available under
the carrier's instrument.  This notion is inherited from carrier
exactification and the foundations framework
\cite{Tsiokos2026CarrierExactification,Tsiokos2026FoundationsII}.  At this
point the carrier is a conceptual object; the concrete \NewLayer{} carriers are
introduced only in \Cref{sec:methods}.

A native strict slot is a strict-realization claim attached to the native state
space of such a carrier.  Its threshold-attachment data are populated at the
\NewLayer{} level, in the sense of \cite{Tsiokos2026FoundationsII}, rather than
imported from the bridge-native \OldLayer{} level.  The realization assertion is
therefore layer-local: it concerns the new-layer description and does not
upgrade the bridge-level slot status.

The native three-slot strict domain is the conjunction of three native strict
slots on the same new-layer carrier: readout-agreement, quadratic-witness, and
transition-witness.  The domain is closed only when all three slots are
strict-realized at that same layer and the status-family discipline of the
audit framework is preserved.  This closed domain is the structural object on
which strict-test reassembly later operates in \Cref{sec:primary-results}.
\Cref{subsec:strict-test-family-descent-compatibility} introduces strict test
families and descent compatibility, \Cref{subsec:non-factorization-baseline}
records the non-factorization boundary, and
\Cref{subsec:old-layer-status-table} records the \OldLayer{} status ledger.

\subsection{Strict test families and descent compatibility}\label{subsec:strict-test-family-descent-compatibility}

A strict test family is a finite family of test predicates over the slots of a
closed strict domain.  It satisfies the inherited test-family discipline:
mutual exclusivity, exhaustivity, readout-category coverage, and family
structure \cite{Tsiokos2026CarrierToEvent}.  It is accepted only after the
family-criteria gate and the required controls have been passed under the
foundations audit discipline \cite{Tsiokos2026FoundationsII}.  In
\PaperNine{}, this is a layer-local notion.  Acceptance at a \NewLayer{} carrier
does not by itself imply acceptance at the bridge-native \OldLayer{} level,
where no accepted strict test family is asserted.

Descent compatibility is the corresponding level-crossing notion.  A descent
map sends new-layer evidence back to the \OldLayer{} description through the
loss-record framework of \cite{Tsiokos2026FoundationsII}.  Compatibility means
that the image of new-layer strict realization is a real-provisional
pair-observable record together with a boundary record on slot identities, not
an old-layer strict-realization claim.  This relation lets new-layer strict
realization coexist with the inherited \OldLayer{} status without upgrading the
old layer or destabilizing the new one.  \Cref{subsec:non-factorization-baseline}
records the irreducibility condition, \Cref{sec:primary-results} supplies the
strict-test audits, and \Cref{sec:secondary-results} supplies the descent and
control-survival audits.

\subsection{Non-factorization boundary}\label{subsec:non-factorization-baseline}

In the foundations sense, a layer transition is a strict extension when the
new-layer carrier cannot be reconstructed from the old-layer description alone
\cite{Tsiokos2026FoundationsII}.  Equivalently, the descent map admits no
section recovering the carrier from its old-layer image.  The boundary between
layers is then non-factorizing: it carries \NewLayer{} structure that does not
factor through any \OldLayer{} presentation.  \PaperNine{} uses this as a
conceptual definition here.  The carrier-side certificate for the exactified
carriers and the preservation of non-factorization under descent are recorded
in
\Cref{subsec:non-factorization-certificate,subsec:non-factorization-preservation}.

\subsection{Old-layer baseline status}\label{subsec:old-layer-status-table}

The inherited \OldLayer{} profile is organized by six dimensions.  Its
pair-observable support is real-provisional in the sense of \PaperEight{}
\cite{Tsiokos2026CarrierToEvent}.  Its bridge-level threshold status is not
met under the inherited threshold-and-witness discipline; only provisional
inspection is admitted there \cite{Tsiokos2026FoundationsII}.  Its control
status is control-ready, because the bridge-level controls validated in
\PaperEight{} remain available.  Its strict-slot status is mixed: the
quadratic-levelset and transition-stability slots are realized as scoped event
positives, but the readout-agreement slot remains unrealized, so the strict
three-slot domain does not close at \OldLayer{}.  Its test-family status is
that no accepted strict test family is asserted at the old layer.  Its descent
status is compatibility target: any higher-layer descent lands in the
\OldLayer{} description through the foundations loss-record discipline.

This prose ledger records the starting state needed for the rest of the paper.
The extended status table appears in \Cref{app:extended-material}.  The status
is inherited and retained throughout \PaperNine{} as the \OldLayerStatus{}, not
established here as a new result.  The non-factorizing character of the boundary
is the conceptual condition recorded in
\Cref{subsec:non-factorization-baseline}.

\subsection{Layer-localization picture}\label{subsec:layer-localization-figure}

\Cref{fig:layer-localization} summarizes the layer-local picture established in
\Cref{subsec:h2-bridge-baseline,subsec:new-layer-carrier-native-slots,subsec:strict-test-family-descent-compatibility,subsec:non-factorization-baseline,subsec:old-layer-status-table}.

\begin{figure}[tbp]
  \centering
  \small
  \setlength{\fboxsep}{6pt}
  \begin{tabular}{@{}c@{\hspace{0.04\textwidth}}c@{\hspace{0.04\textwidth}}c@{}}
    \fbox{\begin{minipage}{0.31\textwidth}
        \textbf{Old layer (\OldLayer{})}
        \begin{itemize}
          \setlength{\itemsep}{0.2em}
          \setlength{\leftmargin}{1.2em}
          \item scoped event positives
          \item real-provisional bridge compatibility
          \item non-factorization boundary
        \end{itemize}
      \end{minipage}
    }
    &
    \begin{minipage}{0.18\textwidth}
      \centering
      \(\Longleftarrow\)\\[0.35em]
      descent\\
      compatibility only\\
      no old-layer overwrite
    \end{minipage}
    &
    \fbox{\begin{minipage}{0.31\textwidth}
        \textbf{New layer (\NewLayer{})}
        \begin{itemize}
          \setlength{\itemsep}{0.2em}
          \setlength{\leftmargin}{1.2em}
          \item exactified carrier
          \item native three-slot strict domain
          \item accepted strict test families
        \end{itemize}
      \end{minipage}
    }
  \end{tabular}
  \caption{Layer-localization picture for \PaperNine{}: bridge-native \OldLayer{} retains scoped event positives and the real-provisional non-factorization boundary; the \NewLayer{} carrier carries the closed native three-slot strict domain and the accepted strict test families; descent maps new-layer evidence to old-layer compatibility without inducing strict realization at \OldLayer{}.}
  \label{fig:layer-localization}
\end{figure}
 
With the conceptual layer separation fixed, \Cref{sec:methods} constructs the
\NewLayer{} carriers to which the later strict-realization audits apply.
 \section{PICA route and carrier exactification}\label{sec:methods}

\subsection{PICA route overview}\label{subsec:pica-route-overview}

The methodology of \PaperNine{} is the PICA route.  PICA abbreviates
``Primitive Interaction Closure Algebra'' in the underlying dynamics
formalism; in the present paper the name denotes the protocol-indexed
carrier-audit route that promotes one-layer-up candidate objects, internalizes
their protocol coordinates, exactifies the resulting carriers, and admits
sources under an explicit audit policy.  The route operates above the
bridge-native \OldLayer{} event level inherited from \PaperEight{}
\cite{Tsiokos2026CarrierToEvent}.  It is therefore not a direct repair of the
old layer, but the method by which \PaperNine{} constructs the native
\NewLayer{} carrier on which later strict-test questions are posed.

The component stages are inherited from the fixed-support sequence.  Promotion
uses the criteria of \cite{Tsiokos2026PromotionCriteria}; the one-layer-up
structure being promoted is constrained by the internal-structure analysis of
\cite{Tsiokos2026InternalStructure}; carrier exactification supplies the native
carrier, transport law, and readout law of
\cite{Tsiokos2026CarrierExactification}; and the source-policy and audit
discipline are those of the foundations framework
\cite{Tsiokos2026FoundationsII}.  The remainder of this section records these
stages in order: carrier formation and source policy
(\Cref{subsec:pica-readiness,subsec:pica-source-policy-profile}), claim-boundary
discipline (\Cref{subsec:pica-claim-boundary-report}), candidate inputs and the
promotion gate (\Cref{subsec:candidate-objects-promotion-gate}), carrier
exactification (\Cref{subsec:carrier-exactification}), native state space with
transport and readout law (\Cref{subsec:native-state-transport-readout}),
non-factorization (\Cref{subsec:non-factorization-certificate}), and the table
of exactified carriers (\Cref{subsec:exactified-carriers-table}).

\subsection{Carrier-formation prerequisites}\label{subsec:pica-readiness}

Before the audits of \Cref{sec:primary-results,sec:secondary-results} are
applied, the \NewLayer{} carriers are already fixed by the upstream construction
stages.  Promotion of one-layer-up candidate objects follows
\cite{Tsiokos2026PromotionCriteria}; the promoted structure is constrained by
the internal-structure analysis of \cite{Tsiokos2026InternalStructure}; and
carrier exactification supplies the exactified carrier objects of
\cite{Tsiokos2026CarrierExactification}.  The end state of this finite audited
construction is the three carriers used below: the lifted-package surface, the
protocol-indexed higher-witness surface, and the completion-family surface.
Their existence is taken as fixed from this point on.  The later sections test
threshold attachment, slot realization, descent, and control survival over
these carriers.  The exactification step that closes the construction is
recorded in \Cref{subsec:carrier-exactification}.

\subsection{Source-policy profile}\label{subsec:pica-source-policy-profile}

The source-policy profile is the foundations-level rule that decides which
source classes may support which role claims under a declared instrument
\cite{Tsiokos2026Foundations,Tsiokos2026FoundationsII}.  It assigns classes
such as committed-state evidence, audited cell records, simulation traces,
ablation records, and fallback evidence to admissibility statuses for roles
such as pair-observable status, threshold realization, and slot realization.
A role projection may be asserted only when the active policy admits the
relevant source class for the role and instrument in question.  The policy is
therefore an admissibility rule, not itself a projection or a realization gate.

\PaperNine{} specializes this rule layer-locally.  The bridge-native policy
inherited from \PaperEight{} remains the policy for the \OldLayer{} description
\cite{Tsiokos2026CarrierToEvent}; the PICA-bound policy instantiated here is
the policy for the native \NewLayer{} carriers.  In particular,
committed-state evidence is admitted for new-layer threshold realization when
the attachment contract of
\Cref{subsec:threshold-attachment-requirements} is satisfied.  The same source
class is not thereby admitted as old-layer strict-realization evidence; at the
bridge-native level it remains within the provisional-inspection discipline of
\PaperEight{}.  Thus the source policy is not a global table applied uniformly
across levels, but a pair of layer-indexed policies tied to the active
instrument and carrier.

The policy guarantees admissibility discipline: no role projection is asserted
without source evidence admitted at the layer where the claim is made, and
admission at one layer does not transfer automatically to another.  The audit
record remains separate from the policy decision.  The policy says which source
classes may count; the audit records that the required evidence is present and
meets the attachment contract.  It does not by itself strict-realize a slot or
accept a test family.  Those decisions are made by the gates of
\Cref{sec:primary-results}, after promotion and exactification have supplied
the carrier inputs of
\Cref{subsec:candidate-objects-promotion-gate,subsec:carrier-exactification}.
Nor does the policy license public-shadow inversion or a canonical lift; those
remain explicit nonclaims taken up again in \Cref{sec:analysis}.

\subsection{Claim-boundary discipline}\label{subsec:pica-claim-boundary-report}

\PaperNine{} uses the claim and nonclaim discipline of the foundations layer
\cite{Tsiokos2026FoundationsII} to separate three kinds of statements.  The
positive claim is the strict-extension result: the construction of native
\NewLayer{} carriers, closure of the three-slot strict domain, acceptance of
the strict test families, and survival of the required controls.  The retained
compatibility claim is the \OldLayer{} real-provisional non-factorization
boundary, including the scoped event positives inherited from \PaperEight{}
\cite{Tsiokos2026CarrierToEvent}.  The explicit nonclaims are stronger
structures not asserted here: an event package, a proto-effect fragment, an
effect algebra, a Born valuation, composition readiness, faithful public
compression, a lawful public carrier, and a canonical selected lift.

This discipline is enforced throughout the paper rather than only at the final
statement.  The gates of \Cref{sec:primary-results} attach the acceptance claim
to the new-layer domain and its control-survival condition; the controls of
\Cref{sec:secondary-results}, especially
\Cref{subsec:protocol-honesty-claim-boundary}, check that the accepted families
respect the nonclaim record; \Cref{sec:final-gate} applies the legality
constraints that select the final verdict; and \Cref{sec:analysis} gives the
layer-localization argument explaining why the positive claim does not imply
the stronger structures just disclaimed.

\subsection{Candidate objects and promotion gate}\label{subsec:candidate-objects-promotion-gate}

The candidate objects for the PICA route are the three live one-layer-up
surfaces retained by the carrier-exactification enumeration
\cite{Tsiokos2026CarrierExactification}.  The lifted-package surface records
the packaged higher-level lift data; the protocol-indexed higher-witness
surface organizes candidate witness data by the protocol coordinates that will
remain in scope; and the completion-family surface records the completion-side
family structure needed for later carrier formation.  These surfaces are also
constrained by the internal-structure analysis of
\cite{Tsiokos2026InternalStructure}.  They are named here as methodological
inputs, not as already exactified carriers or strict-test evidence.

The promotion gate applies the six-criterion test of
\cite{Tsiokos2026PromotionCriteria}.  In prose, the candidate must be
sufficient on the designated observables, induce closure of the update law,
remain stable at the second stage under admissible recoarsening, be admissible
against the matched controls, preserve protocol honesty by avoiding smuggled
claims, and remain legible on at least one designated slice.  Each of the
lifted-package, protocol-indexed higher-witness, and completion-family surfaces
passes all six criteria.  The result is the three promoted candidates that
proceed to carrier exactification in \Cref{subsec:carrier-exactification}.

\subsection{Carrier exactification}\label{subsec:carrier-exactification}

Carrier exactification applies the audit framework of
\cite{Tsiokos2026CarrierExactification} to the promoted candidates of
\Cref{subsec:candidate-objects-promotion-gate}.  The audit asks whether a
candidate survives as a carrier-level object with its own native state space,
transport law, and readout law, rather than collapsing back to lower-level
quotient baselines.  Its criteria require readout sufficiency, coherent
transport and readout laws, non-reducibility against the lower quotient
baselines, honest section behavior, protocol honesty, robustness under stress
and expansion, and support across the declared seed samples.  These criteria
extend the promotion discipline of \cite{Tsiokos2026PromotionCriteria} from
candidate admissibility to carrier formation.

Applied at the flagship-only protocol-indexed scope inherited from
\PaperEight{}'s carrier analysis \cite{Tsiokos2026CarrierToEvent}, each of the
three promoted candidates passes the exactification audit.  The lifted-package
surface, the protocol-indexed higher-witness surface, and the completion-family
surface therefore produce the three exactified carriers used downstream.  This
scope qualifier is part of the claim discipline: the carriers are exactified at
the present flagship-only protocol-indexed scope, not at a broader cross-seed
or fully public-shadow scope.  Their native state spaces, transport laws, and
readout laws are recorded in \Cref{subsec:native-state-transport-readout}.

\subsection{Native state space, transport law, and readout law}\label{subsec:native-state-transport-readout}

Each exactified carrier carries three structural components of its own: a native
state space, a transport law on that state space, and a readout law for the
observable content presented through the carrier's instrument
\cite{Tsiokos2026CarrierExactification,Tsiokos2026FoundationsII}.  The state
space provides the domain on which the carrier's content lives.  The transport
law specifies how that content is moved by the carrier dynamics.  The readout
law specifies which carrier-internal content is observable under the declared
instrument.  These structures are native to the \NewLayer{} carrier; they are
not inherited as bridge-native \OldLayer{} state, transport, or readout data.

The carrier record also contains a public-shadow split and a forward-invariance
record.  The public-shadow split is nonfaithful: the carrier has a public
projection to the bridge-native surface, but the carrier is not reducible to
that projection, and faithful public compression remains an explicit nonclaim
\cite{Tsiokos2026CarrierExactification}.  The forward-invariance record is the
audit datum stating that the transport law is stable under the operations
permitted by the foundations framework.  The slot construction of
\Cref{sec:primary-results} uses the native state space, transport law, and
readout law as the substrate on which the required slots are realized.  The
non-factorization property relating these carriers to the bridge-native
description is recorded in \Cref{subsec:non-factorization-certificate}.

\subsection{Non-factorization certificate}\label{subsec:non-factorization-certificate}

Each exactified carrier carries a carrier-side non-factorization certificate:
its native state space, transport law, and readout law do not factor through the
bridge-native \OldLayer{} description.  Equivalently, the carrier cannot be
reconstructed from old-layer data alone.  This is the carrier-level instance of
the foundations non-factorization condition for a strict extension
\cite{Tsiokos2026FoundationsII}, applied to the exactified carrier structures
of \cite{Tsiokos2026CarrierExactification}.

The certificate is methodological.  It records why the \NewLayer{} carriers are
not reformulations of bridge-native data and supplies the structural
prerequisite for the final strict-extension verdict of \Cref{sec:final-gate}.
The descent-side statement that this non-factorization is preserved when
new-layer evidence is translated back to old-layer compatibility is separate
and is recorded in \Cref{subsec:non-factorization-preservation}.

\subsection{Exactified carriers}\label{subsec:exactified-carriers-table}

\Cref{tab:exactified-carriers} summarizes the three exactified \NewLayer{}
carriers produced by the PICA route, recording their source candidates, native
structure, public-shadow split, carrier-side non-factorization, and scope.

\begin{table}[tbp]
  \centering
  \caption{Exactified \NewLayer{} carriers of \PaperNine{}: source candidate, native structure, public-shadow split, non-factorization, and scope.}
  \label{tab:exactified-carriers}
  \small
  \begin{tabularx}{\textwidth}{@{}>{\raggedright\arraybackslash}p{0.20\textwidth}>{\raggedright\arraybackslash}X>{\raggedright\arraybackslash}p{0.17\textwidth}>{\raggedright\arraybackslash}p{0.16\textwidth}>{\raggedright\arraybackslash}p{0.18\textwidth}@{}}
    \toprule
    Carrier & Native structure & Public-shadow split & Non-factorization & Scope \\
    \midrule
    Lifted-package surface
      & native state, transport, and readout laws
      & nonfaithful projection
      & carrier-side certificate
      & flagship-only protocol-indexed \\
    Protocol-indexed higher-witness surface
      & native state, transport, and readout laws
      & nonfaithful projection
      & carrier-side certificate
      & flagship-only protocol-indexed \\
    Completion-family surface
      & native state, transport, and readout laws
      & nonfaithful projection
      & carrier-side certificate
      & flagship-only protocol-indexed \\
    \bottomrule
  \end{tabularx}
\end{table}
 
With the carrier substrate fixed, \Cref{sec:primary-results} turns to threshold
attachment, slot realization, and strict-test-family construction on these
carriers.
 \section{Native strict domain and strict-test acceptance}\label{sec:primary-results}

\subsection{Threshold attachment requirements}\label{subsec:threshold-attachment-requirements}

The threshold-realization claim used in this section is governed by the
attachment discipline of the scoped exact-six framework
\cite{Tsiokos2026FoundationsII}.  Before any role projection is asserted, the
description must already carry the relevant raw feature cluster, a threshold
annotation attached to that cluster, a witness schema specifying what would
support the claim, a collapse-case record, a level annotation, and an audit
record.  The requirement is preconditional: threshold data and witness data are
records on the description, not labels inferred after a channel has been marked
active.

\PaperNine{} inherits this contract from \PaperEight{}
\cite{Tsiokos2026CarrierToEvent} and specializes it in two ways.  First,
attachment is checked on the native state space of the \NewLayer{} carrier; the
bridge-native \OldLayer{} level cannot by itself supply attachment for
new-layer realization without a recorded bridge.  Second, the admissible source
policy is the PICA-bound policy fixed in \Cref{sec:methods}, replacing the
bridge-native source policy used in \PaperEight{}.  The next subsection applies
this contract; the present subsection fixes only the requirements.

\subsection{Source-backed threshold evidence}\label{subsec:source-backed-threshold-evidence}

The attachment audit is applied to the new-layer carriers produced by the PICA
route of \Cref{sec:methods}: the lifted-package surface, the protocol-indexed
higher-witness surface, and the completion-family surface.  On each carrier's
native state space, the six records required by
\Cref{subsec:threshold-attachment-requirements} are present: the raw feature
cluster, threshold annotation, witness schema, collapse-case record, level
annotation, and audit record are populated before any active role projection is
claimed.  This is the inherited attachment discipline of
\cite{Tsiokos2026FoundationsII}, applied at the \NewLayer{} carrier level rather
than at the bridge-native level inherited from \PaperEight{}
\cite{Tsiokos2026CarrierToEvent}.

The admitted source class for this threshold stage is committed-state evidence
under the PICA-bound source-of-truth policy of \Cref{sec:methods}.  That
admission is layer-specific: the same source class is not admitted as
threshold-realization evidence for the bridge-native \OldLayer{} level, where
the inherited source-policy discipline permits only provisional inspection.  The
new-layer policy therefore supplies source-backed threshold evidence, on each of
the three carriers, for the readout-agreement role.  The gate that decides
whether this evidence strict-realizes the readout slot is
\Cref{subsec:new-layer-readout-slot-gate}; the present subsection records only
that the evidence exists and satisfies the attachment contract.

\subsection{New-layer readout slot gate}\label{subsec:new-layer-readout-slot-gate}

The readout slot gate consumes the source-backed threshold evidence of
\Cref{subsec:source-backed-threshold-evidence} together with the PICA-bound
source-policy contract fixed in \Cref{sec:methods}.  Its output is a
strict-realization decision for the readout-agreement slot at the native
\NewLayer{} carrier level.  This is a layer-local decision in the sense of the
inherited slot-realization discipline \cite{Tsiokos2026FoundationsII}: evidence
attached to the new-layer carrier does not by itself upgrade the bridge-native
\OldLayer{} slot status, and the gate does not infer an active projection from
threshold data without the intervening slot-realization check.

Applied to the three carriers, the gate accepts strict realization of the
readout-agreement slot on each carrier's native state space.  The
bridge-native \OldLayer{} readout slot is unchanged by this acceptance and
remains within the \OldLayerStatus{}.  Strict realization of this one slot does
not close the strict three-slot domain and does not license strict-test
acceptance.  The companion quadratic-witness and transition-witness slots are
treated in
\Cref{subsec:native-quadratic-witness-slot-construction,subsec:native-transition-witness-slot-construction},
and the all-slot closure test is deferred to
\Cref{subsec:all-slot-gate-and-strict-domain-closure}.

\subsection{Native quadratic-witness slot construction}\label{subsec:native-quadratic-witness-slot-construction}

The second required slot is the quadratic-witness slot.  Its source material is
the retained exact-finite quadratic structure isolated in the internal-structure
and carrier-exactification papers: the self-quadratic readout family
(\texttt{q\_readout\_l2\_self}), the gap-type quadratic family
(\texttt{q\_readout\_gap}), and the recoarsening-consistency family
(\texttt{q\_coarsening\_consistency})
\cite{Tsiokos2026InternalStructure,Tsiokos2026CarrierExactification}.  In
\PaperNine{} these families are not imported as bridge-native names.  They are
revalidated on the native state space of the \NewLayer{} carrier introduced in
\Cref{sec:methods}: each family is carried by the new-layer transport and
readout law, and each carries the attachment records required by
\Cref{subsec:threshold-attachment-requirements} under the inherited
threshold-and-witness discipline \cite{Tsiokos2026FoundationsII}.

The resulting quadratic-witness slot is strict-realized at the new layer on
each of the three carriers.  This is a native construction on the exactified
carrier, not inheritance-by-name from the bridge-native level.  The quadratic
event family retained by \PaperEight{} remains scoped to the \OldLayer{} side
and does not by itself license the new-layer slot.  The transition-witness slot
is treated in \Cref{subsec:native-transition-witness-slot-construction}, and
the strict three-slot domain is not closed until
\Cref{subsec:all-slot-gate-and-strict-domain-closure}.

\subsection{Native transition-witness slot construction}\label{subsec:native-transition-witness-slot-construction}

The third required slot is the transition-witness slot.  Its source material is
the transition structure isolated in the holonomy-with-memory paper: the
distinction between the current quotient and the predictive quotient, predictive
witness pairs separating histories that agree at the current level but differ
predictively, and the lifted loop action on the predictive quotient
\cite{Tsiokos2026HolonomyMemory}.  It also uses the carrier-level transition
records supplied by carrier exactification, where transition-stability data is
carried at the refinement level \cite{Tsiokos2026CarrierExactification}.  In
\PaperNine{} these structures are not imported as bridge-native names.  They
are revalidated on the native state space of the \NewLayer{} carrier introduced
in \Cref{sec:methods}: the carrier transport law acts on the predictive
quotient, the predictive witnesses persist on the carrier surface, and the
attachment records required by \Cref{subsec:threshold-attachment-requirements}
are populated for the transition-witness slot under the inherited
threshold-and-witness discipline \cite{Tsiokos2026FoundationsII}.

The resulting transition-witness slot is strict-realized at the new layer on
each of the three carriers.  This is a native transition construction on the
exactified carrier, not inheritance-by-name from the bridge-native level.  The
transition-stability event family retained from \PaperEight{}
\cite{Tsiokos2026CarrierToEvent} remains scoped to the \OldLayer{} side and
does not by itself license the new-layer slot.  The strict three-slot domain is
not closed until \Cref{subsec:all-slot-gate-and-strict-domain-closure}, and no
strict-test acceptance is licensed here.

\subsection{All-slot gate and strict-domain closure}\label{subsec:all-slot-gate-and-strict-domain-closure}

The all-slot gate consumes the three slot-realization results established in
\Cref{subsec:new-layer-readout-slot-gate,subsec:native-quadratic-witness-slot-construction,subsec:native-transition-witness-slot-construction}.
It returns closure of the \NewLayer{} strict three-slot domain only when the
readout-agreement, quadratic-witness, and transition-witness slots are all
strict-realized on the same exactified carrier.  The gate therefore enforces
the status-family discipline of the foundations layer
\cite{Tsiokos2026FoundationsII}: realization at one slot does not imply
realization at another, and the conjunction is checked explicitly.  The
bridge-native \OldLayer{} domain is not closed by this gate, which operates on
new-layer evidence under the layer-local attachment discipline fixed in
\Cref{subsec:threshold-attachment-requirements}.

On each of the three carriers, all three required slots are strict-realized on
the carrier's native state space.  The \NewLayer{} strict three-slot domain is
therefore closed on each carrier.  This closes the slot-realization phase of
\Cref{sec:primary-results}; \Cref{subsec:native-strict-domain-table} records
the closed domain, and \Cref{subsec:strict-test-family-construction} begins the
separate test-family construction.  Domain closure is the structural
prerequisite for that construction, not a test-acceptance result.

\subsection{Native strict three-slot domain}\label{subsec:native-strict-domain-table}

\Cref{tab:native-strict-domain} records the closed \NewLayer{} strict
three-slot domain produced by the all-slot gate: the source structure for each
slot, the carriers on which it is supported, its strict-realization status, and
the layer-local claim scope.

\begin{table}[tbp]
  \centering
  \caption{Native PICA \NewLayer{} strict three-slot domain at \PaperNine{}: source structure, supporting carriers, strict-realization status, and claim scope per slot.}
  \label{tab:native-strict-domain}
  \small
  \begin{tabularx}{\textwidth}{@{}>{\raggedright\arraybackslash}p{0.16\textwidth}>{\raggedright\arraybackslash}X>{\raggedright\arraybackslash}p{0.21\textwidth}>{\raggedright\arraybackslash}p{0.16\textwidth}>{\raggedright\arraybackslash}p{0.18\textwidth}@{}}
    \toprule
    Slot & Source structure & Carriers supporting & Strict realization & Claim scope \\
    \midrule
    Readout-agreement slot
      & PICA-bound source-policy admission of committed-state evidence under the threshold-realization framework of \cite{Tsiokos2026FoundationsII}
      & lifted-package surface; protocol-indexed higher-witness surface; completion-family surface
      & strict-realized at the new layer
      & new-layer slot only; bridge-native \OldLayer{} slot unchanged \\
    Quadratic-witness slot
      & three retained quadratic families of \cite{Tsiokos2026InternalStructure} and \cite{Tsiokos2026CarrierExactification}
      & lifted-package surface; protocol-indexed higher-witness surface; completion-family surface
      & strict-realized at the new layer
      & new-layer slot only; bridge-native \OldLayer{} slot unchanged \\
    Transition-witness slot
      & predictive-quotient and lifted-loop-action structure of \cite{Tsiokos2026HolonomyMemory} together with the carrier-level transition records of \cite{Tsiokos2026CarrierExactification}
      & lifted-package surface; protocol-indexed higher-witness surface; completion-family surface
      & strict-realized at the new layer
      & new-layer slot only; bridge-native \OldLayer{} slot unchanged \\
    \bottomrule
  \end{tabularx}
\end{table}
 
This table completes the slot-domain record used by the test-family
construction of \Cref{subsec:strict-test-family-construction}.

\subsection{Strict test-family construction}\label{subsec:strict-test-family-construction}

The construction step takes the closed \NewLayer{} strict three-slot domain of
\Cref{subsec:all-slot-gate-and-strict-domain-closure}, summarized in
\Cref{tab:native-strict-domain}, and assembles a candidate strict test family on
each carrier.  Each candidate is a finite family of test predicates over the
readout-agreement, quadratic-witness, and transition-witness slots, constructed
with the aim of satisfying the test-family discipline inherited from
\PaperEight{}: mutual exclusivity, exhaustivity, readout-category coverage, and
family structure \cite{Tsiokos2026CarrierToEvent}.  The construction is carried
out under the finite-audited-interaction discipline of
\cite{Tsiokos2026FoundationsII}.  It is native to the carrier: the family uses
the strict-realized slots of that carrier and is not assembled from the
bridge-native event family retained by \PaperEight{}.

There are three constructed candidates, one per supporting carrier: the test
family on the lifted-package surface, the test family on the protocol-indexed
higher-witness surface, and the test family on the completion-family surface.
This is the structural construction step.  The family-criteria gate is treated
in \Cref{subsec:family-criteria-gate}, and the final accepted set is reported
in \Cref{subsec:accepted-strict-test-families}.

\subsection{Test-family admissibility gate}\label{subsec:family-criteria-gate}

The admissibility gate applies the test-family discipline inherited from
\PaperEight{} to the three candidates constructed in
\Cref{subsec:strict-test-family-construction} \cite{Tsiokos2026CarrierToEvent}.
It checks that each candidate is finite and that its predicates form an
internally consistent partition of the closed \NewLayer{} strict three-slot
domain: mutual exclusivity prevents overlap between test alternatives,
exhaustivity is evaluated against the closed domain of
\Cref{subsec:all-slot-gate-and-strict-domain-closure}, readout-category
coverage is evaluated against the readout-agreement slot on the carrier's
native state space, and family structure checks that the predicates remain a
single coherent test family rather than unrelated slotwise assertions.  These
checks remain within the finite-audited-interaction framework of
\cite{Tsiokos2026FoundationsII}.

Each of the three constructed candidates satisfies the gate on its supporting
carrier.  Each meets the finiteness, mutual-exclusion, exhaustivity,
readout-coverage, and family-structure requirements over its native strict
three-slot domain.  This gate passage is still not the acceptance declaration:
the accepted strict test families are named only in
\Cref{subsec:accepted-strict-test-families}.

\subsection{Accepted strict test families}\label{subsec:accepted-strict-test-families}

Acceptance follows the gate passage established in
\Cref{subsec:family-criteria-gate}.  On each supporting carrier, the constructed
family is accepted as a strict test family of the \NewLayer{} strict three-slot
domain: it qualifies under the inherited test-family discipline as a finite,
mutually exclusive, exhaustive, readout-covering, and structurally coherent test
partition over the closed domain \cite{Tsiokos2026CarrierToEvent}.  The
acceptance is local to \PaperNine{}'s native carrier construction and to its
scoped protocol-indexed mode.  It does not extend to the bridge-native
\OldLayer{} level of \PaperEight{}, which retains its scoped event positivity
but no accepted strict test family.

The accepted set comprises \AcceptedTestCount{} families, one for each carrier:
the strict test family on the lifted-package surface, the strict test family on
the protocol-indexed higher-witness surface, and the strict test family on the
completion-family surface.  Acceptance at this stage remains subject to the
control-survival analysis of \Cref{sec:secondary-results}.  The table in
\Cref{subsec:accepted-strict-test-families-table} records the per-family
acceptance attributes; the canonical operational identifiers for the three
families are recorded in \Cref{app:artifact-map}.

\subsection{Accepted strict test families table}\label{subsec:accepted-strict-test-families-table}

\Cref{tab:accepted-strict-test-families} records the three accepted PICA
\NewLayer{} new-layer strict test families, their source carriers, their member
slots, their family-criteria status, and the control-survival condition attached
to the acceptance claim.

\begin{table}[tbp]
  \centering
  \caption{Accepted PICA \NewLayer{} new-layer strict test families on the carriers of \PaperNine{}: source carrier, member slots, criteria status, scope, and acceptance condition.  Test-family identifiers are recorded in \Cref{subsec:accepted-strict-test-families}.}
  \label{tab:accepted-strict-test-families}
  \small
  \begin{tabularx}{\textwidth}{@{}>{\raggedright\arraybackslash}p{0.21\textwidth}>{\raggedright\arraybackslash}X>{\raggedright\arraybackslash}p{0.18\textwidth}>{\raggedright\arraybackslash}p{0.16\textwidth}>{\raggedright\arraybackslash}p{0.18\textwidth}@{}}
    \toprule
    Test family & Member slots & Family criteria & Scope & Acceptance \\
    \midrule
    Family on lifted-package surface
      & readout-agreement, quadratic-witness, transition-witness
      & all five criteria pass
      & new-layer; protocol-indexed
      & accepted, conditional on \Cref{sec:secondary-results} \\
    Family on protocol-indexed higher-witness surface
      & readout-agreement, quadratic-witness, transition-witness
      & all five criteria pass
      & new-layer; protocol-indexed
      & accepted, conditional on \Cref{sec:secondary-results} \\
    Family on completion-family surface
      & readout-agreement, quadratic-witness, transition-witness
      & all five criteria pass
      & new-layer; protocol-indexed
      & accepted, conditional on \Cref{sec:secondary-results} \\
    \bottomrule
  \end{tabularx}
\end{table}
 
With the new-layer strict domain and accepted family record fixed, the analysis
turns in \Cref{sec:secondary-results} to descent and control survival.
 \section{New-to-old descent and control survival}\label{sec:secondary-results}

\subsection{Descent map and loss records}\label{subsec:descent-map}

\Cref{sec:primary-results} establishes strict realization on the native
\NewLayer{} carriers.  The descent map sends those new-layer records back to
the \OldLayer{} description through the loss/lift bookkeeping of the scoped
exact-six framework \cite{Tsiokos2026FoundationsII}.  Its purpose is not to
repeat the strict realization at the old layer.  Rather, it records which
new-layer structure can be translated into old-layer compatibility, which
structure is forgotten in the level crossing, and where the inherited
bridge-native event level of \PaperEight{} remains the relevant target
\cite{Tsiokos2026CarrierToEvent}.  The map is therefore the bridge from
new-layer strict realization to old-layer compatibility statements, not an
upgrade to old-layer strict realization.

Under this descent, the readout-agreement role descends as a real-provisional
pair-observable record, the boundary separating new-layer strict realization
from old-layer compatibility is retained, and the new-layer slot identities are
not collapsed into old-layer slot identities.  The carrier-native state space,
transport law, and quadratic- and transition-witness substructure are forgotten
as old-layer structure; the loss record marks this forgetting as part of the
translation rather than as falsification.  Thus the descent preserves
\OldLayer{} as the \OldLayerStatus{}: compatibility is retained, while
old-layer strict realization is not asserted.  \Cref{subsec:non-factorization-preservation}
records the non-factorization statement explicitly, \Cref{subsec:old-h2-compatibility-status}
records the old-layer compatibility status, and the subsequent control panel
tests the accepted families against comparator constructions.

\subsection{Non-factorization preservation}\label{subsec:non-factorization-preservation}

The descent of \Cref{subsec:descent-map} preserves non-factorization: the
new-layer strict realization established in \Cref{sec:primary-results} does not
factor through any old-layer strict-realization claim.  Equivalently, the
descent map has no section that reconstructs the native \NewLayer{} carrier
from the \OldLayer{} description alone.  This is the formal sense in which the
result is a strict extension in the foundations framework: the layer transition
adds structure that cannot be reduced to the lower-layer description
\cite{Tsiokos2026FoundationsII}.  The inherited bridge-native setting is the
\OldLayer{} level of \PaperEight{} \cite{Tsiokos2026CarrierToEvent}, which
serves as the target of compatibility translation rather than as a source of
new-layer strict realization.

This preservation is what keeps the strict-extension verdict of
\Cref{sec:final-gate} from collapsing into an old-layer claim.  If
factorization were available, the new-layer strict realization could in
principle be read as old-layer strict realization, and the boundary enforced by
the descent would fail.  The structural picture is therefore a one-way
combination: new-layer strict realization descends to old-layer compatibility,
but no inverse reconstruction turns that compatibility into the new-layer
carrier.  \Cref{subsec:old-h2-compatibility-status} records the complementary
old-layer status.

\subsection{Old-layer compatibility status}\label{subsec:old-h2-compatibility-status}

After descent, the bridge-native \OldLayer{} level retains the scoped event
positivity inherited from \PaperEight{} \cite{Tsiokos2026CarrierToEvent}.  The
quadratic level-set event and the transition-stability test remain unchanged,
the readout-agreement slot remains unrealized at the old layer, and no accepted
strict test family is placed there.  The evidence constructed in
\Cref{sec:primary-results} is therefore not imported into \OldLayer{} as strict
realization.  Its old-layer image is the compatibility record described in
\Cref{subsec:descent-map}: a real-provisional pair-observable record on the
readout-agreement role together with a boundary record preventing the new-layer
slot identities from being identified with old-layer slot identities.

This status is part of the strict-extension result, not a negative residue.  It
supplies the boundary condition under which the \NewLayer{} construction is
genuinely new: the \OldLayer{} description is left intact, the
non-factorization statement of \Cref{subsec:non-factorization-preservation} is
preserved, and the descent records compatibility rather than inverse
reconstruction.  The section now turns from this layer-separation picture to
the robustness tests for the accepted families.

\subsection{Null baseline and memory-only controls}\label{subsec:control-panel-null-memory}

The control panel asks whether the accepted families of
\Cref{subsec:accepted-strict-test-families} survive degenerate comparators
rather than merely reflecting the construction of a sufficiently elaborate
carrier.  This is the audited-control role of the foundations framework
\cite{Tsiokos2026FoundationsII}, with control vocabulary inherited from the
fixed-support sequence \cite{Tsiokos2026CarrierToEvent,Tsiokos2026PromotionCriteria}.
The present subsection treats the first two control classes; later subsections
address artifact and boundary-sensitive controls, selected-lift and visibility
perturbations, and protocol-honesty checks.

The null baseline is an uninformative comparator: it preserves the carrier
shape but supplies no candidate-distinguishing structure.  Passing this control
shows that the family is not accepted merely because some carrier was
constructed.  The memory-only comparator supplies history or memory but not the
structure that separates the predictive quotient from the current quotient
\cite{Tsiokos2026HolonomyMemory}.  Passing this control shows that the family
is not accepted merely because memory is present.  The three accepted families
pass both controls on their supporting carriers.  This establishes control
survival only for the two comparator classes treated here; the remaining
classes are handled in
\Cref{subsec:control-panel-artifact-boundary,subsec:selected-lift-visibility-controls,subsec:protocol-honesty-claim-boundary}.

\subsection{Artifact-internalization and boundary controls}\label{subsec:control-panel-artifact-boundary}

The artifact-internalization controls test whether acceptance depends on the
operational bookkeeping by which the carrier was exactified or on the particular
way protocol coordinates were internalized into state.  This control class is
inherited from the promotion and carrier-exactification stages, where candidate
structure must survive variation of the audit record rather than merely
reproduce it \cite{Tsiokos2026PromotionCriteria,Tsiokos2026CarrierExactification}.
Passing the control means that the accepted family does not collapse when the
internalization is varied or when the bookkeeping itself is treated as a
comparator.

The boundary-sensitive control supplies a negative comparator at the edge of
the closed \NewLayer{} strict three-slot domain, in the completion-sensitive
sense inherited from \PaperEight{} \cite{Tsiokos2026CarrierToEvent}.  It asks
whether the family detects only structure internal to its closed domain, rather
than producing spurious positives just outside that boundary.  Both controls
pass for all three accepted families on their supporting carriers.  This
extends the control-survival record but does not complete it; selected-lift and
visibility perturbations are treated in
\Cref{subsec:selected-lift-visibility-controls}, and protocol-honesty and
claim-boundary compliance are treated in
\Cref{subsec:protocol-honesty-claim-boundary}.

\subsection{Selected-lift and visibility-suppression controls}\label{subsec:selected-lift-visibility-controls}

The selected-lift variation control applies the Foundations II discipline that
a level crossing may record selected lifts without making any one lift
canonical \cite{Tsiokos2026FoundationsII}.  The control varies the selected lift
within its admissible alternatives and tests whether the accepted family remains
stable under that variation, using the inherited lift records of the promotion
stage \cite{Tsiokos2026PromotionCriteria}.  Passing the control means that the
family is not parasitic on a particular lift choice; the result persists while
the lift remains noncanonical.

The visibility-suppression control perturbs the instrument-relative record of
what is visible and what is suppressed within the admissible range of the
foundations bookkeeping \cite{Tsiokos2026FoundationsII}.  Passing this control
means that acceptance is not an artifact of a single visibility record or a
single suppression convention.  The same per-family passage holds for
selected-lift variation and visibility-suppression perturbation.  These
controls preserve the public-shadow boundary inherited from \PaperEight{}: the
public compression remains nonfaithful, and no additional public-carrier claim
is licensed
\cite{Tsiokos2026CarrierToEvent}.  The final control class, protocol honesty and
claim-boundary compliance, is treated in
\Cref{subsec:protocol-honesty-claim-boundary}.

\subsection{Protocol honesty and claim-boundary compliance}\label{subsec:protocol-honesty-claim-boundary}

The protocol-honesty control checks that the accepted families remain presented
at the scoped protocol-indexed mode in which they were constructed.  Its role is
to prevent protocol coordinates from being silently elided when the family is
reported as a result, a discipline inherited from the promotion and
carrier-exactification stages
\cite{Tsiokos2026PromotionCriteria,Tsiokos2026CarrierExactification}.  Passing
the control means that the family does not pretend to a public-shadow result
that has not been established; the protocol-indexed scope is preserved.

The claim-boundary control checks the explicit nonclaim record of the
foundations framework \cite{Tsiokos2026FoundationsII}.  The accepted families do
not implicitly construct an event package, a proto-effect fragment, an effect
algebra, a Born valuation, or composition readiness; nor do they assert
faithful public compression, a lawful public carrier, or a canonical selected
lift.  These structures are named here only as disclaimed stronger conclusions.
Passing the control means that the accepted families respect that nonclaim
record rather than carrying stronger structure into the acceptance result.

Each of the three accepted families passes the protocol-honesty and
claim-boundary controls on its supporting carrier, in addition to the controls
of \Cref{subsec:control-panel-null-memory,subsec:control-panel-artifact-boundary,subsec:selected-lift-visibility-controls}.
The control-survival condition attached to
\Cref{subsec:accepted-strict-test-families} and
\Cref{tab:accepted-strict-test-families} is therefore satisfied in full: the
families are accepted on \PaperNine{}'s evidence, with layer separation,
non-factorization preservation, and the inherited claim discipline intact.
\Cref{sec:analysis} gives the corresponding layer-localization and
claim-boundary analysis.
 \section{Claim-boundary analysis and layer localization}\label{sec:analysis}

\subsection{Why strict tests are layer-local}\label{subsec:strict-tests-layer-local}

The strict-test acceptance of \Cref{sec:primary-results} is layer-local because
its supporting conditions are layer-local.  First, the source-policy profile is
not global: the PICA-bound policy of
\Cref{subsec:pica-source-policy-profile}, formalized at schema level in
\Cref{subsec:schema-source-policy-threshold}, admits the relevant source class
for \NewLayer{} threshold realization while the inherited bridge-native policy
of \PaperEight{} does not do so at \OldLayer{}
\cite{Tsiokos2026CarrierToEvent,Tsiokos2026FoundationsII}.  Second, the
threshold-attachment data required by
\Cref{subsec:threshold-attachment-requirements} are populated on the native
state space of the exactified \NewLayer{} carrier, not on the bridge-native
description.  Third, the strict-realization gates of \Cref{sec:primary-results}
operate on that new-layer evidence, and the no-propagation schema of
\Cref{subsec:schema-old-blocks-old-only} prevents their acceptance from being
read as old-layer strict realization.  These conditions make accepted strict
tests an attribute of the new layer's description and audit chain, not of the
layer transition by itself.

The same separation excludes the converse reading.  The inherited \OldLayer{}
status of \PaperEight{} does not propagate as a constraint on \NewLayer{}
test-family acceptance: the unrealized old-layer readout slot blocks old-layer
strict-test closure, but it does not block a new-layer carrier that carries its
own readout slot, its own threshold attachment, and its own native state space.
Likewise, the bridge-native restriction of committed-state evidence to
provisional inspection is replaced, at the new layer only, by the PICA-bound
policy.  The result is a clean layer separation: strict-test acceptance lives at
\NewLayer{}, the inherited real-provisional boundary lives at \OldLayer{}, and
the descent map between them is a compatibility translation rather than an
inheritance of strict realization.

\subsection{Bounded scope: no event package, no shared-event identity}\label{subsec:tests-do-not-imply-event-package}

Strict-test acceptance certifies that, on each \NewLayer{} carrier, the family
criteria pass and the controls survive.  It does not provide a packaging
operator over the accepted families, an identification of shared events across
families, or any compositional structure linking the three carriers as a single
algebraic object.  Under the claim/nonclaim discipline of
\cite{Tsiokos2026FoundationsII}, those are additional structural moves, not
consequences of test-family acceptance.  The event-packaging structures
introduced in \cite{Tsiokos2026LocallyBoolean} and bounded at the event bridge
in \cite{Tsiokos2026CarrierToEvent} therefore remain explicit nonclaims in
\PaperNine{}.

An event package would require a packaging operator that respects
shared-event identification across the three accepted families and remains
audited under the foundations framework.  A shared-event identity claim would
require a verified equivalence among events from different families on a common
comparison base \cite{Tsiokos2026LocallyBoolean}.  Neither structure is
constructed here.  The claim-boundary control of
\Cref{subsec:protocol-honesty-claim-boundary} is the audit-side check that
these structures are not imported into the acceptance result.

\subsection{Higher structures remain future work}\label{subsec:higher-structures-blocked}

The accepted strict-test families are partitions of the closed three-slot
domain at each carrier.  Building a proto-effect fragment would require
identifying effect-shaped structure in those families that survives the
foundations claim discipline; building an effect algebra would require
composition-respecting operations on such effects across the three carriers; a
Born valuation would require a state-functional pairing whose values respect
the effect-algebra structure; and composition readiness would require the
algebraic and state-side constructions to combine coherently.  Each step is a
separate structural construction beyond the strict-test acceptance reported in
\PaperNine{}.  Under \cite{Tsiokos2026FoundationsII}, absence of those
constructions is a nonclaim boundary, not an implied consequence of the
accepted tests.

The blockers are constructive rather than refutative.  These higher structures
remain possible targets for later work; the present paper neither constructs
them nor rules them out.  The claim-boundary discipline of
\Cref{subsec:pica-claim-boundary-report} keeps them explicit as nonclaims, and
the claim-boundary control of \Cref{subsec:protocol-honesty-claim-boundary}
verifies on the audit side that the accepted families respect that boundary.

\subsection{Public-shadow nonfaithfulness as a structural prerequisite}\label{subsec:public-shadow-nonfaithfulness}

The public-shadow split inherited from
\cite{Tsiokos2026CarrierExactification} is part of each \NewLayer{} carrier's
structure: the carrier projects to a bridge-native public surface, but it is
not reducible to that projection.  This nonfaithfulness matters directly for
\PaperNine{}'s strict-extension claim.  If the public projection were faithful,
then the new-layer carrier would be reducible to old-layer public-shadow data,
which would make inverse reconstruction available and would collapse the
non-factorization boundary of
\Cref{subsec:non-factorization-baseline,subsec:non-factorization-certificate}.
Nonfaithfulness is therefore a structural prerequisite of the strict-extension
verdict, not a separate caveat.  The visibility-suppression control of
\Cref{subsec:selected-lift-visibility-controls} verifies on the audit side that
the result is preserved across admissible perturbations of the visibility
record without licensing faithful public compression.

\subsection{Selected-lift nonuniqueness}\label{subsec:selected-lift-nonuniqueness-analysis}

The selected-lift discipline of \cite{Tsiokos2026FoundationsII} allows level
crossings to record selected lifts without making any one lift canonical.
\PaperNine{}'s strict-extension verdict respects this discipline: it does not
claim that a particular lift is uniquely admissible, and the selected-lift
variation control of \Cref{subsec:selected-lift-visibility-controls} certifies
that the accepted families survive admissible variation.  A canonical selected
lift would assert that one section across the layer transition is privileged;
such a claim would conflict with the descent-and-no-section structure of
\Cref{subsec:non-factorization-baseline,subsec:non-factorization-preservation}.
Non-factorization and selected-lift nonuniqueness are therefore complementary:
the \NewLayer{} carrier admits no inverse reconstruction, and no particular
lift across the transition is canonical.  The acceptance result lives within
both constraints.  These constraints shape the discussion of
\Cref{sec:discussion}.
 \section{Formal support and Lean schemas}\label{sec:formal-framework}

\subsection{Theorem schemas overview}\label{subsec:theorem-schemas-overview}

This section records the schema-level support behind the gate logic of
\PaperNine{}.  The formal layer consists of theorem schemas that codify the
discipline by which threshold attachment, slot realization, test-family
acceptance, descent compatibility, and the legality constraints of the final
gate are organized.  The schemas are inherited from the foundations layer
\cite{Tsiokos2026FoundationsII} and follow the convention of \PaperEight{}
\cite{Tsiokos2026CarrierToEvent}: formal support is a typed mirror of the gate
discipline, not an independent derivation of empirical content.

The schemas check legality and gate compatibility.  They do not prove that the
\NewLayer{} carriers exist, that source-backed threshold evidence is present,
or that the controls pass; those empirical premises are supplied by the methods
and results sections \Cref{sec:methods,sec:primary-results,sec:secondary-results}.
The role of this section is conditional: if the empirical premises hold, then
the resulting claims obey the paper's internal discipline.  The schemas proceed
in order: source-policy admission required for threshold realization
(\Cref{subsec:schema-source-policy-threshold}), promotion versus carrier
exactification (\Cref{subsec:schema-promotion-not-exactification}),
exactified-carrier attachment licensing new-layer slots
(\Cref{subsec:schema-attachment-licenses-slot}), new-layer slots licensing
new-layer test-family reassembly (\Cref{subsec:schema-slots-license-tests}),
the real-provisional old layer blocking old-layer tests but not new-layer tests
(\Cref{subsec:schema-old-blocks-old-only}), unrealized required new-layer slots
blocking accepted new-layer strict tests
(\Cref{subsec:schema-unrealized-blocks-test}), and the boundary on what the
formal layer does not prove (\Cref{subsec:schema-legality-not-existence}).

\subsection{Source-policy admission for threshold realization}\label{subsec:schema-source-policy-threshold}

The first schema states that a threshold-realization claim at a layer is
licensed only when the active source-of-truth policy at that layer admits the
source class supplying the threshold evidence for the relevant role
\cite{Tsiokos2026FoundationsII}.  Source-policy admission is therefore a
structural precondition of every active threshold-realization claim.  In
\PaperNine{}, this schema underwrites the layer-locality of the threshold claims
used in \Cref{sec:primary-results}: committed-state evidence is admitted for
threshold realization at the \NewLayer{} carrier under the PICA-bound policy of
\Cref{subsec:pica-source-policy-profile}, while the same source class remains
only provisionally inspectable at the bridge-native \OldLayer{} level under the
inherited policy of \PaperEight{} \cite{Tsiokos2026CarrierToEvent}.

\subsection{Promotion and carrier exactification as distinct stages}\label{subsec:schema-promotion-not-exactification}

The second schema states that promotion of a candidate object, in the sense of
\cite{Tsiokos2026PromotionCriteria}, is distinct from carrier exactification, in
the sense of \cite{Tsiokos2026CarrierExactification}.  A promoted candidate is
not by itself an exactified carrier: carrier exactification requires
already-promoted candidates as input and applies its own audit.  Promotion is
therefore necessary but not sufficient for carrier-level standing.  In
\PaperNine{}, this schema underwrites the staged method of \Cref{sec:methods},
where \Cref{subsec:candidate-objects-promotion-gate} applies the promotion gate
and \Cref{subsec:carrier-exactification} applies the exactification audit with
a separate criterion set.  It also enforces the claim discipline that the three
\NewLayer{} carriers are exactified, not merely promoted, before any
slot-realization claim is asserted in \Cref{sec:primary-results}.

\subsection{Exactified-carrier attachment licenses a native strict slot}\label{subsec:schema-attachment-licenses-slot}

The third schema states a conditional slot-licensing principle.  If an
exactified carrier carries threshold-attachment data on its native state space,
in the sense of \cite{Tsiokos2026FoundationsII}, and the active source policy
admits the supplying source class for threshold realization, then a native
strict slot is licensed at that carrier under the inherited slot-realization
framework.  Equivalently, carrier exactification
\cite{Tsiokos2026CarrierExactification}, threshold attachment, and source-policy
admission are jointly sufficient for new-layer slot licensing.  This schema
underwrites the readout-slot gate of
\Cref{subsec:new-layer-readout-slot-gate} and the native companion-slot
constructions of
\Cref{subsec:native-quadratic-witness-slot-construction,subsec:native-transition-witness-slot-construction}.

\subsection{Strict-domain closure licenses test-family reassembly}\label{subsec:schema-slots-license-tests}

The fourth schema states that closure of the \NewLayer{} strict three-slot
domain at a carrier licenses strict-test-family reassembly at that carrier.
Closure means that the readout-agreement, quadratic-witness, and
transition-witness slots are all strict-realized at the same carrier.  It is
the necessary licensing input to family construction, but acceptance still
requires the family-criteria gate: mutual exclusivity, exhaustivity,
readout-category coverage, and family structure
\cite{Tsiokos2026CarrierToEvent}, under the audit framework of
\cite{Tsiokos2026FoundationsII}.  This schema underwrites the closure result of
\Cref{subsec:all-slot-gate-and-strict-domain-closure} and the construction and
criteria-gate steps of
\Cref{subsec:strict-test-family-construction,subsec:family-criteria-gate}.  It
does not by itself license acceptance; that requires gate passage and the
control panel of \Cref{sec:secondary-results}.

\subsection{Old-layer real-provisional status blocks old-layer tests only}\label{subsec:schema-old-blocks-old-only}

The fifth schema is a no-propagation statement.  Retaining \OldLayer{} at
real-provisional non-factorization status blocks acceptance of strict test
families at \OldLayer{}, but it does not block acceptance at \NewLayer{} when
the evidence is constructed at the new layer under the PICA-bound source policy.
Equivalently, an inherited unrealized old-layer slot constrains old-layer
test-family acceptance only; it is not a constraint on new-layer test-family
acceptance whose threshold evidence, slot realization, and controls are
layer-local.  This is the formal counterpart of the layer-locality framework of
\cite{Tsiokos2026FoundationsII} applied to the inherited setting of
\cite{Tsiokos2026CarrierToEvent}.  It underwrites the principal positive case
of \Cref{sec:final-gate} and the descent-side compatibility statements of
\Cref{subsec:non-factorization-preservation,subsec:old-h2-compatibility-status}.

\subsection{Unrealized required slot blocks new-layer test acceptance}\label{subsec:schema-unrealized-blocks-test}

The sixth schema is the corresponding new-layer blocking statement.  If any
required slot of the \NewLayer{} strict three-slot domain is not strict-realized
at a carrier, then the all-slot gate of
\Cref{subsec:all-slot-gate-and-strict-domain-closure} fails for that carrier,
and no accepted strict test family is licensed there.  Equivalently, the
readout-agreement, quadratic-witness, and transition-witness slots are jointly
necessary for new-layer test-family acceptance; partial slot realization is
insufficient.  This schema is the new-layer counterpart of
\Cref{subsec:schema-old-blocks-old-only}: partial realization at the new layer
does not propagate to test-family acceptance at that layer.  It formalizes the
discipline applied in the empirical run of \PaperNine{}, where the companion
slots had to be constructed before the strict-domain rerun audit could close
and acceptance could become available, in accordance with the audit and
test-family frameworks of \cite{Tsiokos2026FoundationsII} and
\cite{Tsiokos2026CarrierToEvent}.

\subsection{Boundary: schema legality, not empirical existence}\label{subsec:schema-legality-not-existence}

The schemas of this section establish schema legality and gate constraints:
they record the conditions under which \PaperNine{}'s claim discipline is
internally consistent, not what is empirically the case.  They do not prove that
the \NewLayer{} carriers exist, that source-backed threshold evidence is
present, that the controls pass, or that descent compatibility holds.  Those
are empirical statements established by the artifact pipeline described in
\Cref{sec:methods,sec:primary-results,sec:secondary-results}.  This boundary is
the one inherited from \cite{Tsiokos2026FoundationsII} and used throughout the
fixed-support series
\cite{Tsiokos2026CarrierToEvent,Tsiokos2026CarrierExactification,Tsiokos2026PromotionCriteria,Tsiokos2026InternalStructure}.
The appendix \Cref{app:formalization} records the corresponding Lean-side
schema entries and proof structure where formalized; the present section is the
manuscript-side mirror of those entries.
 \section{Final gate}\label{sec:final-gate}

\subsection{Final verdict statement}\label{subsec:final-verdict-statement}

Relative to the event-admissibility endpoint inherited from \PaperEight{}
\cite{Tsiokos2026CarrierToEvent}, \PaperNine{} establishes a strict extension
on fixed support.  More precisely, the PICA route described in
\Cref{sec:methods}, consisting here of promotion, exactification, and
source-policy-admitted carrier construction, yields a \NewLayer{} carrier with a
native three-slot strict domain.  The readout, quadratic-witness, and
transition-witness slots are strict-realized at that layer, and
\AcceptedTestCount{} accepted PICA \NewLayer{} strict test families survive the
required controls.

The old layer does not become a strict-test layer in the present analysis.  It
remains the bridge-native event level of \PaperEight{}, retained throughout as a
\OldLayerStatus{}; non-factorization is preserved.  This separates new-layer
strict realization from old-layer compatibility.  \Cref{subsec:decision-logic}
explains the ordered case selection, and \Cref{subsec:final-claim-boundary}
records the boundary between what the paper claims, what it retains as
compatibility, and what it explicitly disclaims.

\subsection{Decision logic}\label{subsec:decision-logic}

The final gate is an ordered case analysis over the upstream audits from
\Cref{sec:methods,sec:primary-results,sec:secondary-results}.  The principal
positive case is reached when the native \NewLayer{} strict domain closes, the
strict test families on the new-layer carrier survive controls, and descent
preserves old-layer compatibility without turning \OldLayer{} into a
strict-realized layer.  Three weaker positive cases are reserved for incomplete
new-layer evidence: a threshold-only case in which the new-layer readout slot is
strict-realized but test-family reassembly is unavailable; a real-provisional
case in which the route preserves or improves support without a completed
threshold-and-slot gate; and a below-threshold or pending case.

The same case analysis contains five blocking cases and two terminal negative
cases.  A blocker case obtains if promotion supplies no surviving object, if
promotion survives but carrier exactification does not, if exactified threshold
evidence cannot descend compatibly to the old layer, if source-policy admission
blocks the source use, or if public-shadow safety blocks the route.  A
falsification case is legal only when neither old-layer nor new-layer threshold,
readout, real-provisional, or below-threshold evidence remains.  A separate
implementation-failure case is reserved for malformed or unavailable artifacts,
not for mathematical absence.

These constraints are needed because partial evidence can otherwise be read in
more than one way: for example, incomplete new-layer evidence could be mistaken
for falsification if retained provisional evidence were ignored.  The final gate
therefore closes such overlaps.  The falsification case is unavailable in the
presence of retained provisional or below-threshold evidence; the threshold-only
case requires a strict-realized new-layer readout slot; and the strict-extension
case requires the accepted strict test families on the new-layer carrier
together with control survival.  The persistence of the old-layer
real-provisional non-factorization boundary is not a negative result; it is the
old-layer condition under which the new-layer strict extension is licensed.
Artifact provenance for these audits is collected in \Cref{app:artifact-map}.

\subsection{Final claim boundary}\label{subsec:final-claim-boundary}

\Cref{tab:final-claim-boundary} records the final claim boundary: what the paper
claims, what it retains as compatibility, and what it explicitly disclaims.  The
table distinguishes the accepted strict tests on the new-layer carrier, the
retained \OldLayer{} compatibility boundary, and the stronger structures that
remain explicit nonclaims.

\begin{table}[tbp]
  \centering
  \caption{Final claim boundary for \PaperNine{}: results claimed at the new layer, old-layer compatibility boundaries, explicit nonclaims, and supporting artifact families.}
  \label{tab:final-claim-boundary}
  \small
  \begin{tabularx}{\textwidth}{@{}>{\raggedright\arraybackslash}X>{\raggedright\arraybackslash}X>{\raggedright\arraybackslash}p{0.16\textwidth}>{\raggedright\arraybackslash}p{0.15\textwidth}@{}}
    \toprule
    Item & Status & Claimed? & Supporting artifact family \\
    \midrule
    Accepted PICA \NewLayer{} new-layer strict tests
      & accepted; control-surviving
      & yes (count = \AcceptedTestCount)
      & H, I, J, K \\
    \OldLayer{} strict tests
      & not constructed
      & no
      & C, K \\
    \OldLayer{} strict readout realization
      & not realized
      & no
      & C, G, K \\
    \OldLayer{} real-provisional boundary
      & \OldLayerStatus{}
      & yes (as boundary)
      & B, C, G, K \\
    Shared-event identity
      & not assumed
      & no
      & G, I, J, K \\
    Event package
      & not constructed
      & no
      & I, J, K \\
    Proto-effect fragment
      & not constructed
      & no
      & I, J, K \\
    Effect algebra
      & not constructed
      & no
      & I, J, K \\
    Born valuation
      & not constructed
      & no
      & I, J, K \\
    Composition readiness
      & not investigated
      & no
      & I, J, K \\
    Faithful public compression
      & not assumed
      & no
      & G, J, K \\
    Lawful public carrier
      & not claimed
      & no
      & G, J, K \\
    Canonical selected lift
      & not assumed
      & no
      & G, J, K \\
    \bottomrule
  \end{tabularx}
\end{table}
 
The table records the discipline that distinguishes the strict-extension result
from an old-layer strict realization and from any stronger structure not claimed
in this paper.
 \section{Discussion and limitations}\label{sec:discussion}

\subsection{Discussion of the strict-extension result}\label{subsec:discussion-strict-extension}

\PaperNine{} establishes strict-test acceptance as a layer-local result.  The
accepted families are attached to the \NewLayer{} carriers constructed by the
PICA route, while the bridge-native \OldLayer{} inherited from \PaperEight{}
\cite{Tsiokos2026CarrierToEvent} is retained as a real-provisional
non-factorization boundary.  The principal contribution is therefore not a
repair of the old layer, but the layer-localization picture: strict tests are
admissible at the new layer, the inherited old-layer setting is preserved, and
the one-way descent between them records compatibility without inverse
reconstruction.  This is the strict-extension form anticipated by the
foundations framework \cite{Tsiokos2026FoundationsII}.

The result also clarifies the role of controls.  Control survival does not
enlarge the claim beyond the new-layer strict domain; it confirms that the
accepted families are stable under the comparator, perturbation, protocol, and
claim-boundary tests of \Cref{sec:secondary-results}.  The analysis of
\Cref{sec:analysis} is consequently a localization argument: the accepted
families are genuine at \NewLayer{}, compatible with \OldLayer{}, and bounded
against stronger readings.

\subsection{Scope, limitations, and future work}\label{subsec:discussion-scope-limitations}

The main limitations are structural.  The accepted strict test families do not
constitute an event package, a proto-effect fragment, an effect algebra, a Born
valuation, or a composition-ready structure.  Those remain explicit nonclaims,
for the reasons given in
\Cref{subsec:tests-do-not-imply-event-package,subsec:higher-structures-blocked}:
each would require additional construction beyond strict-test acceptance.
They are future-work targets rather than negative results.

Public-shadow nonfaithfulness and selected-lift nonuniqueness also remain part
of the result's scope.  They are not incidental caveats; they are structural
conditions that keep the strict extension from collapsing into a public-shadow
or canonical-lift reconstruction, as discussed in
\Cref{subsec:public-shadow-nonfaithfulness,subsec:selected-lift-nonuniqueness-analysis}.
Finally, \PaperNine{} is scoped to the flagship-only protocol-indexed setting.
Global, cross-seed, and full-core test claims are out of scope here.  Whether
the carrier exactification framework of \cite{Tsiokos2026CarrierExactification}
can be extended to broader support is a methodological question for a successor
paper.
 \section{Conclusion}\label{sec:conclusion}

\PaperNine{} establishes the layer-localization of strict-test admissibility on
fixed support.  The \NewLayer{} carriers constructed by the PICA route admit a
closed three-slot strict domain with three accepted strict test families
surviving controls.  The bridge-native \OldLayer{} of \PaperEight{}
\cite{Tsiokos2026CarrierToEvent} is retained intact as a real-provisional
non-factorization boundary, in the sense of the foundations framework
\cite{Tsiokos2026FoundationsII}.  Descent from new-layer evidence to old-layer
compatibility preserves non-factorization without inducing strict realization
at the old layer.  The result therefore supplies a strict extension rather than
a repair of the inherited bridge-native level.  It does not extend to
event-package construction, shared-event identity, proto-effect or
effect-algebra structure, Born valuation, composition readiness, faithful
public compression, lawful public carriers, or canonical selected lifts.  Those
structures are nonclaims of the present paper and remain future-work targets as
discussed in \Cref{sec:discussion}.
 
\appendix
\section{Definitions and gate predicates}\label{app:definitions}

\subsection{Definitions}\label{subsec:definitions-list}

\begin{description}
\item[\OldLayer{}.]
The bridge-native event level inherited from \PaperEight{}
\cite{Tsiokos2026CarrierToEvent}.  It carries scoped event positives and an
unrealized readout-agreement slot.

\item[\NewLayer{}.]
The new-layer carrier substrate constructed by the PICA route of
\Cref{sec:methods}.  It is the layer at which \PaperNine{} states native strict
slot realization.

\item[Real-provisional non-factorization boundary.]
The status retained for \OldLayer{} throughout \PaperNine{}.  It is neither
falsification nor old-layer strict realization; it is the boundary condition
under which the new-layer strict extension is assessed.

\item[Exactified carrier.]
A carrier-level object surviving the seven-criterion exactification audit of
\cite{Tsiokos2026CarrierExactification}.  Such a carrier has its own native
state space, transport law, and readout law.

\item[Native strict slot.]
A strict-realization claim attached to the native state space of an exactified
carrier.  The threshold-attachment data of
\cite{Tsiokos2026FoundationsII} must be populated at that layer.

\item[Native three-slot strict domain.]
The conjunction of the readout-agreement, quadratic-witness, and
transition-witness native strict slots on the same exactified carrier.  The
domain is closed only when all three slots are strict-realized at that carrier.

\item[Strict test family.]
A finite family of test predicates over the slots of a closed strict domain.
It satisfies the test-family discipline of
\cite{Tsiokos2026CarrierToEvent}, including mutual exclusivity, exhaustivity,
readout-category coverage, and family structure.

\item[Descent compatibility.]
The property that new-layer evidence descends to old-layer compatibility
without inducing old-layer strict realization.  In \PaperNine{}, the descent
image is a real-provisional pair-observable record together with a boundary
record on slot identities.

\item[Non-factorization.]
The property that the new-layer carrier cannot be reconstructed from the
old-layer description alone.  Equivalently, in the sense of
\cite{Tsiokos2026FoundationsII}, the descent map admits no section recovering
the carrier from its old-layer image.
\end{description}

\subsection{Gate predicates}\label{subsec:gate-predicates-list}

\begin{description}
\item[Promotion gate.]
The six-criterion test of \cite{Tsiokos2026PromotionCriteria} applied to
one-layer-up candidate objects in
\Cref{subsec:candidate-objects-promotion-gate}.

\item[Exactification gate.]
The seven-criterion test of \cite{Tsiokos2026CarrierExactification} applied to
promoted candidates in \Cref{subsec:carrier-exactification}.

\item[Threshold gate.]
The source-policy admission gate that decides whether threshold-realization
evidence at a carrier is admissible at that layer, as used in
\Cref{subsec:source-backed-threshold-evidence,subsec:new-layer-readout-slot-gate}.

\item[Slot-realization gate.]
The per-slot gate that decides strict realization at a carrier in
\Cref{subsec:new-layer-readout-slot-gate,subsec:native-quadratic-witness-slot-construction,subsec:native-transition-witness-slot-construction}.

\item[All-slot gate.]
The conjunction-checking gate that closes the strict three-slot domain in
\Cref{subsec:all-slot-gate-and-strict-domain-closure}.

\item[Descent audit.]
The audit that records which structure descends to old-layer compatibility and
which structure is forgotten, as used in \Cref{subsec:descent-map}.

\item[Test-family admissibility gate.]
The family-criteria gate of \Cref{subsec:family-criteria-gate}, which applies
the inherited test-family discipline of \cite{Tsiokos2026CarrierToEvent}.

\item[Control panel.]
The panel of comparator and perturbation tests recorded in
\Cref{subsec:control-panel-null-memory,subsec:control-panel-artifact-boundary,subsec:selected-lift-visibility-controls,subsec:protocol-honesty-claim-boundary}.

\item[Final gate.]
The ordered case selection of \Cref{sec:final-gate}.
\end{description}
 \section{Methods, artifacts, and reproducibility contract}\label{app:methods}

\subsection{Artifact discipline}\label{subsec:artifact-discipline}

The empirical content of \PaperNine{} is recorded in machine-readable artifacts
under stable namespaces, following the convention used through the
fixed-support series
\cite{Tsiokos2026CarrierToEvent,Tsiokos2026CarrierExactification}.  The
PICA-bound construction uses a dedicated namespace, separate from the earlier
route-discovery and bridge-diagnosis namespaces, so that \NewLayer{} carrier
formation, threshold evidence, slot realization, descent, and controls are
tracked without overwriting inherited \OldLayer{} records.  The per-family
catalog is given in \Cref{app:artifact-map}.

The artifact pipeline is finite and audited.  Each stage produces a named
output, and the audit chain checks that downstream stages consume the upstream
outputs on which they depend, with no silent rewriting of the recorded inputs.
End-state asset readiness is therefore a structural condition: every claim in
the paper is supported by an artifact whose existence and content are fixed
before the audit gates of
\Cref{sec:primary-results,sec:secondary-results,sec:final-gate} are applied.

\subsection{Artifact families}\label{subsec:artifact-families}

The artifact pipeline is organized by stage, in the convention of
\cite{Tsiokos2026CarrierToEvent,Tsiokos2026CarrierExactification}.  The
source-policy and exactification artifacts correspond to \Cref{sec:methods};
the threshold and slot-realization artifacts correspond to
\Cref{sec:primary-results}; the descent and control artifacts correspond to
\Cref{sec:secondary-results}; and the final-gate artifact corresponds to
\Cref{sec:final-gate}.  Each family is documented in \Cref{app:artifact-map}
with stage-keyed identifiers and dependency relations.  The body sections refer
to the families by descriptive name; the catalog supplies the per-artifact
identifiers.

\subsection{Reproducibility contract}\label{subsec:reproducibility-contract}

The artifact pipeline of \PaperNine{} follows a standard reproducibility
contract: each derived artifact carries enough metadata to reconstruct its
computation from the upstream artifacts on which it depends, in the spirit of
reproducible computational research \cite{Donoho2010,Sandve2013}.  This is the
same contract used through the fixed-support series
\cite{Tsiokos2026CarrierToEvent,Tsiokos2026CarrierExactification}, adapted here
to the PICA-bound construction and its \NewLayer{} audit chain.

The contract covers the structural artifacts of the audit chain, including
source-policy contracts, exactified-carrier records, threshold and
slot-realization audits, descent records, control-panel summaries, and the
final-gate decision.  Stage-specific re-runs and independent verification are
within scope.  Exhaustive numerical reproducibility of stochastic auxiliary
computations is not within scope by design, because the strict-test-family
acceptance asserted in the paper is structural rather than statistical.  This
scope follows the audit-record discipline of
\cite{Tsiokos2026FoundationsII}.
 \section{Artifact source map}\label{app:artifact-map}

This appendix records the stage-keyed artifact catalog for \PaperNine{}.  The
body sections refer to artifact families by scientific role; the catalog below
gives the corresponding machine-readable source files.  The convention follows
the fixed-support artifact discipline of
\cite{Tsiokos2026CarrierToEvent,Tsiokos2026CarrierExactification} and the
audit-record framework of \cite{Tsiokos2026FoundationsII}.

\begin{description}
\item[Inherited route-discovery.]
Part A records the inherited route atlas, candidate placements, comparison
matrices, and route-level gate outputs that precede the \OldLayer{} diagnosis.
The stage is represented by
\begin{itemize}
\item \path{results/derived/ticket33a_route_atlas.json},
\path{results/derived/ticket33a_route_comparison_matrix.json},
\path{results/derived/ticket33a_route_hypothesis_registry.json};
\item \path{results/derived/ticket33a_candidate_readout_agreement_registry.json},
\path{results/derived/ticket33a_candidate_test_registry.json},
\path{results/derived/ticket33a_missing_slot_placement_atlas.json};
\item \path{results/derived/ticket33a_domain_audit_summary.json},
\path{results/derived/ticket33a_final_route_gate.json},
\path{results/derived/ticket33a_partB_recommendation.json};
\item \path{results/derived/ticket33a_claim_boundary_manifest.json},
\path{results/derived/ticket33a_artifact_contract_report.json},
\path{results/derived/ticket33a_asset_index.json}.
\end{itemize}

\item[Inherited \OldLayer{} diagnosis.]
Part B records the bridge-native diagnosis inherited from \PaperEight{} and
used in \Cref{sec:background,sec:problem-setup}.  The stage is represented by
\begin{itemize}
\item \path{results/derived/ticket33b_bridge_candidate_registry.json},
\path{results/derived/ticket33b_bridge_source_inventory.json},
\path{results/derived/ticket33b_bridge_threshold_recovery_summary.json};
\item \path{results/derived/ticket33b_readout_agreement_gate.json},
\path{results/derived/ticket33b_threshold_status_decision.json},
\path{results/derived/ticket33b_strict_slot_realization_audit.json};
\item \path{results/derived/ticket33b_strict_test_reassembly_summary.json},
\path{results/derived/ticket33b_test_gate_summary.json},
\path{results/derived/ticket33b_final_gate_decision.json};
\item \path{results/derived/ticket33b_inherited_nonclaims.json},
\path{results/derived/ticket33b_control_bundle_summary.json},
\path{results/derived/ticket33b_asset_index.json}.
\end{itemize}

\item[Pre-PICA diagnostic Part C.]
The early Part C records collect the diagnostic material that identifies the
one-layer-up route later isolated as the PICA construction.  The stage is
represented by
\begin{itemize}
\item \path{results/derived/ticket33c_foundations_alignment_contract.json},
\path{results/derived/ticket33c_import_rule_registry.json},
\path{results/derived/ticket33c_import_rule_summary.json};
\item \path{results/derived/ticket33c_one_layer_up_candidate_registry.json},
\path{results/derived/ticket33c_new_layer_formation_certificate.json},
\path{results/derived/ticket33c_nonfactorization_obstruction_report.json};
\item \path{results/derived/ticket33c_source_policy_contract.json},
\path{results/derived/ticket33c_source_policy_validation_summary.json},
\path{results/derived/ticket33c_threshold_candidate_registry.json},
\path{results/derived/ticket33c_threshold_source_audit_summary.json};
\item \path{results/derived/ticket33c_partC_gate_invariants.json},
\path{results/derived/ticket33c_partC_pica_reset_manifest.json},
\path{results/derived/ticket33c_partC_restart_readiness.json}.
\end{itemize}

\item[PICA reset and source policy.]
This stage records the dedicated PICA namespace, source-policy profile, warning
resolution, and imported bridge material used by \Cref{sec:methods}.  The stage
is represented by
\begin{itemize}
\item \path{results/derived/ticket33c_pica_namespace_acceptance_certificate.json},
\path{results/derived/ticket33c_pica_source_policy_contract.json},
\path{results/derived/ticket33c_pica_source_policy_validation_summary.json};
\item \path{results/derived/ticket33c_pica_warning_taxonomy.json},
\path{results/derived/ticket33c_pica_warning_resolution_ledger.json},
\path{results/derived/ticket33c_pica_claim_boundary_report.json};
\item \path{results/derived/ticket33c_pica_bridge_adapter_manifest.json},
\path{results/derived/ticket33c_pica_bridge_import_diff.json},
\path{results/derived/ticket33c_pica_bridge_source_index.json};
\item \path{results/derived/ticket33c_pica_bundle_input_freeze.json},
\path{results/derived/ticket33c_pica_pipeline_run_manifest.json},
\path{results/derived/ticket33c_pica_artifact_contract_report.json}.
\end{itemize}

\item[PICA carrier formation and exactification.]
This stage records the promoted candidates, exactified carriers, native carrier
laws, and carrier-side non-factorization certificates consumed by
\Cref{sec:methods}.  The stage is represented by
\begin{itemize}
\item \path{results/derived/ticket33c_pica_promotion_candidate_registry.json},
\path{results/derived/ticket33c_pica_promotion_gate_summary.json},
\path{results/derived/ticket33c_clean_new_layer_formation_certificate.json};
\item \path{results/derived/ticket33c_pica_exactification_gate_summary.json},
\path{results/derived/ticket33c_pica_exactified_carrier_registry.json},
\path{results/derived/ticket33c_operator_exactification_registry_pica.json};
\item \path{results/derived/ticket33c_pica_lawful_shell_summary.json},
\path{results/derived/ticket33c_pica_forward_invariance_registry.json},
\path{results/derived/ticket33c_pica_law_readout_parity_gate.json};
\item \path{results/derived/ticket33c_clean_h6_new_layer_readiness_decision.json},
\path{results/derived/ticket33c_readout_law_summary_pica.json},
\path{results/derived/ticket33c_transport_law_summary_pica.json}.
\end{itemize}

\item[PICA threshold, slot realization, descent, and strict-domain audits.]
This stage records the threshold-evidence audits, slot-realization gates,
new-to-old descent records, and strict-domain closure data used by
\Cref{sec:primary-results,sec:secondary-results}.  The stage is represented by
\begin{itemize}
\item \path{results/derived/ticket33c_pica_new_layer_threshold_candidate_registry.json},
\path{results/derived/ticket33c_pica_new_layer_threshold_source_audit_summary.json},
\path{results/derived/ticket33c_pica_new_layer_slot_realization_gate.json},
\path{results/derived/ticket33c_pica_new_layer_slot_status_registry.json};
\item \path{results/derived/ticket33c_pica_new_layer_quadratic_witness_registry.json},
\path{results/derived/ticket33c_pica_new_layer_quadratic_slot_summary.json},
\path{results/derived/ticket33c_pica_new_layer_transition_witness_registry.json},
\path{results/derived/ticket33c_pica_new_layer_transition_slot_summary.json};
\item \path{results/derived/ticket33c_pica_all_slot_threshold_candidate_registry.json},
\path{results/derived/ticket33c_pica_companion_slot_threshold_audit_summary.json},
\path{results/derived/ticket33c_pica_all_slot_realization_gate.json},
\path{results/derived/ticket33c_pica_all_slot_status_registry.json};
\item \path{results/derived/ticket33c_pica_new_to_old_descent_audit_summary.json},
\path{results/derived/ticket33c_pica_old_layer_compatibility_registry.json},
\path{results/derived/ticket33c_pica_strict_domain_rerun_audit.json},
\path{results/derived/ticket33c_pica_strict_domain_status_registry.json}.
\end{itemize}

\item[PICA accepted tests, control panel, and final gate.]
This stage records the accepted strict test families, the control panel, and
the final claim-boundary decision used by
\Cref{sec:primary-results,sec:secondary-results,sec:final-gate}.  The stage is
represented by
\begin{itemize}
\item \path{results/derived/ticket33c_pica_strict_test_reassembly_summary.json},
\path{results/derived/ticket33c_pica_accepted_test_family_registry.json},
\path{results/derived/ticket33c_pica_control_panel_summary.json};
\item \path{results/derived/ticket33c_pica_final_gate_decision.json},
\path{results/derived/ticket33c_pica_findings_digest.json},
\path{results/derived/ticket33c_pica_asset_index.json};
\item \path{results/derived/ticket33c_pica_final_implementation_readiness.json},
\path{results/derived/ticket33c_pica_final_handoff_cleanup_report.json},
\path{results/derived/ticket33c_pica_final_terminology_contract.json}.
\end{itemize}
\end{description}

All listed artifacts live under \path{results/derived/} in the repository.
Their dependencies and rerun expectations are governed by the reproducibility
contract of \Cref{subsec:reproducibility-contract}.
 \section{Lean formalization notes}\label{app:formalization}

\subsection{Lean theorem schemas}\label{subsec:lean-theorem-schemas}

The schemas of \Cref{sec:formal-framework} are mirrored on the Lean side as
typed entries in the Mathlib-style formalization convention used throughout the
fixed-support series
\cite{Tsiokos2026FoundationsII,Tsiokos2026CarrierExactification,Tsiokos2026PromotionCriteria,Tsiokos2026CarrierToEvent}.
Each manuscript schema has a corresponding typed-record entry, and together
these entries give the Lean view of \PaperNine{}'s claim discipline.  The
entries record preconditions and schema claims: source-policy admission appears
as a typed predicate, exactified-carrier attachment as a structural condition,
and strict-domain closure as a conjunction of the three \NewLayer{} slot
statuses.  They do not contain constructive proofs of the empirical premises;
they record the implication structure under which the manuscript's claim
discipline is internally consistent.

\subsection{Conditional proof structure}\label{subsec:lean-proof-structure}

Each Lean-side schema is supported by a typed proof under explicit assumptions
matching the schema's preconditions.  The proofs are conditional: given the
inputs named by the schema, such as source-policy admission, exactified-carrier
attachment, or slot strict-realization, the claimed implication holds in the
Lean type theory.  This follows the inherited formalization convention of
\cite{Tsiokos2026FoundationsII} and the carrier-side developments of
\cite{Tsiokos2026CarrierExactification,Tsiokos2026PromotionCriteria}.  The
proofs do not derive the empirical premises themselves.  The Lean side does not
construct a particular \NewLayer{} carrier for \PaperNine{} or compute that
committed-state evidence is admitted by the active policy; those are statements
of the artifact pipeline, not of the type theory.  The development is therefore
proof-search-tractable but empirically underdetermined by design: it records the
discipline, not the data.

\subsection{Schema legality and empirical content}\label{subsec:lean-boundary}

The Lean development summarized in
\Cref{subsec:lean-theorem-schemas,subsec:lean-proof-structure} establishes
schema legality and conditional proof structure; it does not establish
empirical witness existence.  This is the same boundary stated in
\Cref{subsec:schema-legality-not-existence}: the appendix records the Lean-side
mirror of the manuscript schemas, not a stronger claim.  The empirical content
is supplied by the artifact pipeline of
\Cref{sec:methods,sec:primary-results,sec:secondary-results}.  This boundary is
consistent with the convention used throughout the fixed-support series
\cite{Tsiokos2026FoundationsII,Tsiokos2026CarrierToEvent,Tsiokos2026CarrierExactification,Tsiokos2026PromotionCriteria,Tsiokos2026InternalStructure}.
 \section{Extended tables and nonclaims ledger}\label{app:extended-material}

\subsection{Extended \OldLayer{} status across audit dimensions}\label{subsec:extended-table-h2-status}

\Cref{tab:extended-h2-status} expands the inherited \OldLayer{} status
summarized in \Cref{subsec:old-layer-status-table}.

\begin{table}[tbp]
\centering
\footnotesize
\caption{Extended status ledger for \OldLayer{} at the close of \PaperEight{}:
status across the six audit dimensions, source of inheritance, and the
disposition retained throughout \PaperNine{}.}
\label{tab:extended-h2-status}
\begin{tabularx}{\textwidth}{@{}>{\raggedright\arraybackslash}p{0.16\textwidth}>{\raggedright\arraybackslash}X>{\raggedright\arraybackslash}p{0.18\textwidth}>{\raggedright\arraybackslash}X@{}}
\toprule
Dimension & Status at \OldLayer{} & Inherited from & Carried forward as \\
\midrule
Pair-observable support
  & Real-provisional support for the readout-agreement role.
  & \PaperEight{} \cite{Tsiokos2026CarrierToEvent}.
  & Old-layer compatibility record. \\
Threshold status
  & Not met; provisional inspection only under the inherited threshold-and-witness discipline.
  & \cite{Tsiokos2026FoundationsII,Tsiokos2026CarrierToEvent}.
  & Reason the readout slot is not strict-realized at \OldLayer{}. \\
Control status
  & Control-ready at the bridge-native level.
  & \PaperEight{} \cite{Tsiokos2026CarrierToEvent}.
  & Inherited control availability; not old-layer test acceptance. \\
Strict-slot realization
  & Two of three slots realized as scoped event positives; readout-agreement slot unrealized.
  & \PaperEight{} \cite{Tsiokos2026CarrierToEvent}.
  & Baseline gap addressed by the \NewLayer{} construction. \\
Test-family acceptance
  & No accepted strict test family asserted at \OldLayer{}.
  & \PaperEight{} \cite{Tsiokos2026CarrierToEvent}.
  & Layer-local boundary; new-layer acceptance does not descend. \\
Descent compatibility
  & Compatibility target for higher-layer descent under loss-record discipline.
  & \cite{Tsiokos2026FoundationsII}, applied to \PaperEight{} \cite{Tsiokos2026CarrierToEvent}.
  & Descent target preserving non-factorization. \\
\bottomrule
\end{tabularx}
\end{table}

\subsection{Extended control survival}\label{subsec:extended-table-control-survival}

\Cref{tab:extended-control-survival} expands the control panel reported in
\Cref{subsec:control-panel-null-memory,subsec:control-panel-artifact-boundary,subsec:selected-lift-visibility-controls,subsec:protocol-honesty-claim-boundary}.
The three accepted families pass each control class uniformly.

\begin{table}[tbp]
\centering
\small
\caption{Extended control-survival table.  Each of the three accepted strict
test families on its supporting carrier passes each control class.}
\label{tab:extended-control-survival}
\begin{tabular}{@{}lll@{}}
\toprule
Control class & Outcome (per family) & Verdict \\
\midrule
Null baseline & Passes & Control survived. \\
Memory-only & Passes & Control survived. \\
Artifact-internalization & Passes & Control survived. \\
Boundary-sensitive negative & Passes & Control survived. \\
Selected-lift variation & Passes & Control survived. \\
Visibility/suppression perturbation & Passes & Control survived. \\
Protocol honesty and claim-boundary & Passes & Control survived; accepted. \\
\bottomrule
\end{tabular}
\end{table}

\subsection{Complete nonclaims ledger}\label{subsec:nonclaims-ledger}

This ledger collects the structures that \PaperNine{} explicitly does not
assert, under the claim/nonclaim discipline of
\cite{Tsiokos2026FoundationsII}.

\begin{description}
\item[Old-layer strict test closure.]
This would assert a closed strict test-family domain at \OldLayer{}.  It remains
a nonclaim because the bridge-native readout-agreement slot is not
strict-realized.

\item[Old-layer readout strict realization.]
This would assert strict realization of the readout-agreement role at
\OldLayer{}.  It remains a nonclaim because committed-state evidence is not
admitted for that strict role under the inherited bridge-native policy of
\PaperEight{}.

\item[Shared-event identity.]
This would identify the new-layer accepted families with old-layer event
objects.  It remains a nonclaim because descent records compatibility and
forgetting, not identity of events across layers.

\item[Event package.]
This would assert an event-algebraic package derived from the accepted families
with additional packaging structure.  It remains a nonclaim because the paper
constructs accepted strict test families only, not an event package.

\item[Proto-effect fragment.]
This would assert a fragment of effect-like structure on the accepted families.
It remains a nonclaim because no such fragment is constructed or audited here.

\item[Effect algebra.]
This would assert a full algebraic effect structure.  It remains a nonclaim
because the strict-test acceptance result does not supply the operations,
identities, or laws required for an effect algebra.

\item[Born valuation.]
This would assert a valuation rule of Born type on the accepted families.  It
remains a nonclaim because no valuation map or valuation law is constructed.

\item[Composition readiness.]
This would assert that the accepted families are prepared for composition with
other structures in the series.  It remains a nonclaim because the paper proves
strict-test acceptance and control survival, not a composition interface.

\item[Faithful public compression.]
This would assert a faithful public projection of the new-layer carrier.  It
remains a nonclaim because the public-shadow split is explicitly non-faithful.

\item[Lawful public carrier.]
This would assert a public carrier satisfying the relevant carrier laws.  It
remains a nonclaim because the exactified carriers are protocol-indexed
\NewLayer{} carriers, not lawful public carriers.

\item[Canonical selected lift.]
This would assert a distinguished selected lift independent of admissible
variation.  It remains a nonclaim because selected-lift variation is controlled
without choosing a canonical lift.

\item[Global, cross-seed, or full-core tests.]
This would assert that the accepted strict test families extend beyond the
scoped protocol-indexed setting of \PaperNine{}.  It remains a nonclaim because
the paper's evidence is confined to the fixed-support and scoped carrier
construction used in the audit.
\end{description}

\subsection{Extended new-to-old descent audit}\label{subsec:extended-table-descent-audit}

\Cref{tab:extended-descent-audit} expands the descent and compatibility
statements of
\Cref{subsec:descent-map,subsec:non-factorization-preservation,subsec:old-h2-compatibility-status}.
The descent profile is uniform across the three carriers; the table records the
shared profile.

\begin{table}[tbp]
\centering
\small
\caption{Extended new-to-old descent audit.  The profile is uniform across the
lifted-package surface, the protocol-indexed higher-witness surface, and the
completion-family surface.}
\label{tab:extended-descent-audit}
\begin{tabularx}{\textwidth}{@{}>{\bfseries\raggedright\arraybackslash}p{0.28\textwidth}X@{}}
\toprule
Field & Value \\
\midrule
Carriers
  & Lifted-package surface; protocol-indexed higher-witness surface; completion-family surface. \\
Retained data
  & Readout-agreement role descends as a real-provisional pair-observable, with a boundary record on slot identities. \\
Forgotten data
  & Carrier-native state space, transport law, quadratic-witness substructure, and transition-witness substructure. \\
Non-factorization preserved
  & Yes; the descent map admits no inverse reconstruction of the carrier from its old-layer image. \\
Old-layer status after descent
  & \OldLayerStatus{}; readout slot remains unrealized at \OldLayer{}. \\
\bottomrule
\end{tabularx}
\end{table}

The table is the appendix form of the prose ledger in
\Cref{subsec:old-layer-status-table}; it records inherited status, not a new
result of \PaperNine{}.

\subsection{Extended new-layer threshold witness audit}\label{subsec:extended-table-threshold-audit}

\Cref{tab:extended-threshold-audit} expands the per-carrier threshold evidence
reported in \Cref{subsec:source-backed-threshold-evidence}.  As with the
descent audit, the per-carrier profile is uniform; the table records the
shared profile.

\begin{table}[tbp]
\centering
\small
\caption{Extended threshold-witness audit.  The profile is uniform across the
three exactified carriers.}
\label{tab:extended-threshold-audit}
\begin{tabularx}{\textwidth}{@{}>{\bfseries\raggedright\arraybackslash}p{0.32\textwidth}X@{}}
\toprule
Field & Value \\
\midrule
Carriers
  & Lifted-package surface; protocol-indexed higher-witness surface; completion-family surface. \\
Source class admitted
  & Committed-state evidence under the PICA-bound source-of-truth policy of \Cref{subsec:pica-source-policy-profile}. \\
Six attachment fields populated
  & Raw feature cluster, threshold annotation, witness schema, collapse-case record, level annotation, and audit record are all populated on the carrier's native state space. \\
Threshold attachment audit
  & Passes on each carrier. \\
Result
  & Source-backed threshold evidence at \NewLayer{} for the readout-agreement role on each of the three carriers. \\
\bottomrule
\end{tabularx}
\end{table}
 
\clearpage
\bibliographystyle{plainnat}
\bibliography{references}

\end{document}
